% !TEX root = ../../mixed-sklyanin-cubics.tex
% Current result: 2026-11-22 / v62.
\section{因子間関係の PBW 剛性と二次境界の混合変形}
\label{ifd:sec:interfactor}

前節までの比較は因子別の表示を用いて構成した。
本節では、固定 tensor 境界と厳密な中心変数を持つ全二次 PBW 表示が、
元の生成元のままその範囲に入ることを証明する。
その結果、指定中心三次を持つ全てのそのような表示に、
既存の正規形、smoothness 判定と可換 Ku 比較を適用できる。
次に、最高次の二次境界も変える次数零変形を調べる。
中心三次の一次持上げ障害は Euler 混合類を排除するが、
Clifford 因子の中心平方を使う混合類は残る。
一つの ordered pair に対応する九方向については、
全係数上の平坦な二次代数族と正則中心三次商を構成する。
その周囲代数は全係数で相対 Koszul かつ非捩れ $6$-Calabi--Yau である。
中心商には分離した原点を含む相対 smooth proper な開集合が存在する。
中心変数を一つ加えた五次元型に加え、rotation の形式認識を使う四次元型でも、
原点近傍の可換 Ku 実現と \'etale 局所的な実比較核を得る。

% !TEX root = ../../mixed-sklyanin-cubics.tex
% Current result: 2026-11-22 / v61.
\subsection{固定 tensor 境界に対する PBW 関係の因子別分解}
\label{ifp:subsec:factorization}

本小節では、因子間の可換性を変形後の仮定から外す。
固定した標準 Sklyanin 因子の tensor 積を最高 active 次数の境界に持ち、
中心変数の作用を厳密に保つ二次 PBW 表示では、
因子間の全補正と、内部関係に現れる他因子の生成元が零になる。
結論は元の生成元による関係式の等式であり、形式的な同型類だけの主張ではない。

PBW 性を三次の関係へ還元する一般原理は
Braverman--Gaitsgory \cite[Main Theorem~0.5, Lemma~3.3]{CMOBravermanGaitsgory}
による。以下の必要条件の証明では、その交換子恒等式を直接に次数分解する。
Gaddis \cite[Proposition~5.1]{CMOGaddisPBW} は essentially regular な代数の
tensor 積が再び essentially regular となることを示しているが、
最高次が tensor 積である全 PBW 変形の分解を主張する結果ではない。
より近い先行結果として Walton--Witherspoon
\cite[Hypothesis~2.6, Proposition~6.2]{CMOWaltonWitherspoonBraided} は
$(S\otimes S^{\mathrm{op}})\#\C C_2$ の定数型 PBW 変形を分類している。
そこでは scalar 値の因子間補正は零であり、二捩れの場合に現れる追加の混合補正は
非自明な群元に値を持つ。その群元は active 生成元と可換でない。
本小節では任意の固定因子の組、線形補正、任意個の厳密な中心変数、
非被約係数環も扱う。同文献の Problem~7.2 が提起する一般の braided smash product
の線形補正を分類するものではない。この範囲の分解定理自体の初出は未確認である。

\subsubsection{二つの低次数単射と PBW 係数の取り出し}

$k\ge2$ とし、滑らかな標準点曲線を持つ固定因子
$S_i=S(A_i,B_i,C_i)$ を取る。各因子の三次元生成空間を
$V_i=\C x_i\oplus\C y_i\oplus\C z_i$ とし、関係空間を
$\mathcal R_i\subset V_i\otimes V_i$ と書く。標準の三関係は
\begin{equation}\label{ifp:eq:sklyanin-relations}
\begin{aligned}
 r_{i,x}&=A_i y_i\otimes z_i+B_i z_i\otimes y_i+C_i x_i\otimes x_i,\\
 r_{i,y}&=A_i z_i\otimes x_i+B_i x_i\otimes z_i+C_i y_i\otimes y_i,\\
 r_{i,z}&=A_i x_i\otimes y_i+B_i y_i\otimes x_i+C_i z_i\otimes z_i.
\end{aligned}
\end{equation}
ここで $A_iB_iC_i\ne0$ である。
以下の二つの単射には、そのうち $C_i\ne0$ だけで十分である。

\begin{lemma}[負一次導分と一次中心の消滅]
\label{ifp:lem:low-degree}
各 $S_i$ は
\[
       \operatorname{Der}_{-1}(S_i)=0,
             \qquad Z(S_i)_1=0
\]
を満たす。より具体的には、次の二つの有限次元 $\C$ 線型写像は単射である。
\begin{equation}\label{ifp:eq:two-injections}
\begin{aligned}
 d_i&:V_i^*\longrightarrow\operatorname{Hom}_{\C}(\mathcal R_i,V_i),\\
 d_i(\lambda)\left(\sum_a v_a\otimes v'_a\right)
   &=\sum_a\bigl(\lambda(v_a)v'_a+\lambda(v'_a)v_a\bigr),\\
 \zeta_i&:V_i\longrightarrow\operatorname{Hom}_{\C}(V_i,(S_i)_2),\\
 \zeta_i(q)(v)&=[v,q].
\end{aligned}
\end{equation}
従って任意の可換 $\C$-代数 $R$ で tensor した写像も単射である。
\end{lemma}

\begin{proof}
次数 $-1$ の導分は $V_i\to\C$ で決まり、関係を保つ条件は
$d_i(\lambda)=0$ である。例えば
\[
 d_i(\lambda)(r_{i,x})
 =2C_i\lambda(x_i)x_i
 +(A_i+B_i)\bigl(\lambda(y_i)z_i+\lambda(z_i)y_i\bigr).
\]
この式の $x_i$ 係数と巡回した二式の対応する係数から
$\lambda(x_i)=\lambda(y_i)=\lambda(z_i)=0$ を得る。
特に、これら三係数を $2C_i$ で割る写像が $d_i$ の左逆を与える。

次に、式\eqref{ifp:eq:sklyanin-relations}の各関係は互いに異なる平方の
非零係数を持つので
\[
                  \mathcal R_i\cap\bigwedge^2V_i=0.
\]
$q\in V_i$ が一次中心元なら、各 $v\in V_i$ について
$v\otimes q-q\otimes v$ はこの共通部分に属するので零である。
三次元空間の全ての $v$ に対してこの等式を満たす $q$ は零である。
これで $\zeta_i$ も単射となる。
有限次元 $\C$ 線型単射は左逆を持つため、どちらも任意の $R$ への tensor 後に単射である。
被約性や閉点での判定は使わない。
\end{proof}

$V=\bigoplus_iV_i$、$B_0=\bigotimes_iS_i$ と置く。
$i<j$ に対し $\mathcal R_{ij}$ を
$v\otimes w-w\otimes v$、$v\in V_i,w\in V_j$ で張られる空間とすれば、
\[
 \mathcal R=\bigoplus_i\mathcal R_i
               \oplus\bigoplus_{i<j}\mathcal R_{ij},
             \qquad B_0=T_{\C}(V)/(\mathcal R).
\]
任意の可換 $\C$-代数 $R$、$m\ge1$、
$P=R[\mathbf u]=R[u_1,\ldots,u_m]$ を取る。
$V_R=V\otimes_{\C}R$ と書き、$P_d$ は中心変数に関する斉次次数 $d$ の部分とする。
中心変数は最初から $P$ の係数として扱い、全生成元と厳密に可換とする。
全二次補正とは、任意の $R$ 線型写像
\[
 \ell:\mathcal R\otimes R\longrightarrow P_1\otimes_RV_R,
       \qquad b:\mathcal R\otimes R\longrightarrow P_2
\]
による表示
\begin{equation}\label{ifp:eq:full-presentation}
 B=T_P(V\otimes_{\C}P)
        \big/\bigl(r-\ell(r)-b(r):r\in\mathcal R\bigr)
\end{equation}
をいう。ここでは内部関係の $\ell$ に全ての因子の生成元を許し、
因子間の交換関係にも $\ell$ と $b$ の双方を許す。
全生成元と $u_s$ の内部次数は一であり、$R$ の内部次数は零である。
一方、active 濾過 $F$ では $V$ の次数を一、$P$ の次数を零とする。
PBW という仮定は、最高 active 次数からの自然な全射が同型
\begin{equation}\label{ifp:eq:pbw}
                 B_0\otimes_{\C}P\xrightarrow{\ \sim\ }\operatorname{gr}_F B
\end{equation}
となることである。

\begin{lemma}[中心次数の係数の独立性]
\label{ifp:lem:coefficient-extraction}
式\eqref{ifp:eq:pbw}の下で、$B_0$ の固定した斉次語基底の像は $B$ の $P$ 自由基底になる。
$I=(u_1,\ldots,u_m)\subset P$ とすると、自然な写像は代数同型
\begin{equation}\label{ifp:eq:center-associated-graded}
       \operatorname{gr}_I B
       \simeq (B_0\otimes_{\C}R)\otimes_R\operatorname{Sym}_R(R^m)
\end{equation}
を与える。右辺の $R^m$ の次数は中心イデアル次数一である。
これらの同定は任意の係数変更 $R\to R'$ と両立する。
\end{lemma}

\begin{proof}
斉次語基底の任意の元を同じ語で $B$ へ持ち上げる。
最高 active 次数を順に消すことで、それらは $B$ を $P$ 上で張る。
有限な一次関係があれば、その最高 active 次数を
式\eqref{ifp:eq:pbw}へ送って最高次係数を消し、帰納的に全係数が零となる。
従って表示した語は $P$ 自由基底である。
各 $I^qB$ はその基底の $I^qP$ 係数部分なので、中心イデアルの随伴加群は右辺に一致する。
積の中心最低次項は、式\eqref{ifp:eq:full-presentation}で
$\ell,b$ を消した $B_0$ の積であるから、同定は代数同型である。
自由基底と表示の恒等式を tensor すれば、非平坦な係数変更についても同じ結論を得る。
\end{proof}

\subsubsection{混合係数の厳密な消滅}

\begin{theorem}[固定 tensor 境界の PBW 因子分解]
\label{ifp:thm:factorization}
式\eqref{ifp:eq:full-presentation}が PBW 条件\eqref{ifp:eq:pbw}を満たすとする。
このとき元の生成元において
\begin{equation}\label{ifp:eq:vanishing}
 \ell(\mathcal R_{ij})=0,\qquad b(\mathcal R_{ij})=0\quad(i<j),
       \qquad \ell(\mathcal R_i)\subset P_1\otimes_R(V_i\otimes R)
\end{equation}
である。従って各 $i$ に対し
\[
 B_i=T_P(V_i\otimes_{\C}P)
       \big/\bigl(r-\ell(r)-b(r):r\in\mathcal R_i\bigr)
\]
と置くと、各 $B_i$ は $S_i\otimes_{\C}P$ の PBW 変形であり、
指定した生成元を保つ $P$ 代数同型
\begin{equation}\label{ifp:eq:tensor-factorization}
                    B\simeq B_1\otimes_P\cdots\otimes_PB_k
\end{equation}
がある。逆に、各因子の PBW 変形のこの tensor 積は
式\eqref{ifp:eq:pbw}を満たす。
結論に平行移動や形式的な座標変更は不要であり、任意の係数変更とも両立する。
中心変数が零個の場合は、同じ結論が $B=B_0\otimes R$ として自明に成り立つ。
\end{theorem}

\begin{proof}
\textbf{中心一次の交換子恒等式.}
$\ell(r)=\sum_su_s\ell_s(r)$ と書く。
$i\ne j$、$w\in V_j$、$r=\sum_av_a\otimes v'_a\in\mathcal R_i$ を固定する。
$B$ の中で実際の交換子を取ると
\[
 [w,r]=\sum_a\bigl([w,v_a]v'_a+v_a[w,v'_a]\bigr)
                 =[w,\ell(r)],
\]
となる。最後の等式では $b(r)\in P$ の中心性を用いた。
この恒等式の $I/I^2$ における $u_s$ 係数を取り、
補題\ref{ifp:lem:coefficient-extraction}を使うと
\begin{equation}\label{ifp:eq:first-commutator}
 \sum_a\bigl(\ell_s([w,v_a])v'_a
                  +v_a\ell_s([w,v'_a])\bigr)
                  =[w,\ell_s(r)]
          \quad\text{in }(B_0)_2\otimes R
\end{equation}
を得る。ここで $[w,v]$ は関係空間の元 $w\otimes v-v\otimes w$ をも表す。
これは三次の重なりの中心一次成分である。

\textbf{因子間の線形補正.}
$h\ne i$ を取り、$\ell_s([w,v])$ の $V_h$ 成分を
\[
                  \sum_\nu\lambda_\nu(v)e_{h,\nu},
           \qquad \lambda_\nu:V_i\otimes R\longrightarrow R
\]
と書く。ただし $e_{h,\nu}$ は $V_h$ の基底である。
式\eqref{ifp:eq:first-commutator}を $B_0$ の因子別多重次数で分解する。
右辺で非零となり得るのは $V_j$ 成分との交換子だけなので、その次数は $2\mathbf e_j$ である。
一方、左辺の次数 $\mathbf e_i+\mathbf e_h$ 成分は
\[
                  \sum_\nu d_i(\lambda_\nu)(r)\otimes e_{h,\nu}
                         \quad\text{in }V_i\otimes V_h\otimes R
\]
である。$i\ne j$、$h\ne i$ より両者の次数は異なり、この成分は零である。
全ての $r\in\mathcal R_i$ に対して成立するので、
補題\ref{ifp:lem:low-degree}の $d_i\otimes R$ の単射性から全ての $\lambda_\nu$ が零となる。
従って $\ell_s([w,v])\in V_i\otimes R$ である。
$i,j$ を交換して同じ議論を行うと、この元は $V_j\otimes R$ にも属する。
両空間の共通部分は零なので $\ell_s(\mathcal R_{ij})=0$ を得る。
これは全ての $s,i,j$ について成り立つ。

\textbf{内部関係の他因子出力.}
前段の結論を式\eqref{ifp:eq:first-commutator}へ戻すと
$[w,\ell_s(r)]=0$ となる。
$\ell_s(r)$ の $V_j$ 成分を $q_j$ とすれば、
他の因子の成分は $w$ と可換なので $[w,q_j]=0$ である。
$w\in V_j$ は任意であり、$\zeta_j\otimes R$ の単射性から $q_j=0$ を得る。
全ての $j\ne i$ について適用し、$\ell_s(\mathcal R_i)\subset V_i\otimes R$ を得る。

\textbf{中心二次の因子間補正.}
ここまでで混合交換子は厳密に $[w,v]=b([w,v])\in P_2$ となった。
同じ内部関係 $r\in\mathcal R_i$ との交換子は
\begin{equation}\label{ifp:eq:second-commutator}
 \sum_a\bigl(b([w,v_a])v'_a+v_ab([w,v'_a])\bigr)
               =\sum_su_s[w,\ell_s(r)]\in I^3B
\end{equation}
を満たす。右辺が $I^3B$ に入るのは
$\ell_s(r)\in V_i\otimes R$ と $[V_j,V_i]\subset P_2$ による。
$b([w,v])=\sum_{s\le t}u_su_t\beta_{st}(v)$ と展開して
式\eqref{ifp:eq:second-commutator}を $I^2B/I^3B$ に移すと、
各 $u_su_t$ の係数は $d_i(\beta_{st})(r)=0$ である。
再び $d_i\otimes R$ の単射性から $\beta_{st}=0$ となり、
$b(\mathcal R_{ij})=0$ を得る。
この段階でも二次の混合単項式 $u_su_t$ を個別に扱っており、
中心方向への点別特殊化を用いていない。

\textbf{表示と PBW 基底の分解.}
式\eqref{ifp:eq:vanishing}により、$B$ の関係は各 $B_i$ の内部関係と
異なる因子の生成元の可換関係だけになる。
tensor 積の普遍性から式\eqref{ifp:eq:tensor-factorization}が従う。
各 $B_i$ の PBW 性も確かめる。
自然な写像の合成
\[
 S_i\otimes_{\C}P\longrightarrow\operatorname{gr}_F B_i
          \longrightarrow\operatorname{gr}_F B
                 \simeq B_0\otimes_{\C}P
\]
は $i$ 番目の因子の包含であり、単射である。
最初の写像は最高次関係からの全射なので同型となる。
従って $B_i$ は PBW で、$B_i\to B$ も濾過の最高次を見れば単射である。
逆方向は各因子の $P$ 自由 PBW 基底を tensor すれば従う。
全ての消滅は $R$ 上の等式として証明され、用いた単射は分裂した $\C$ 線型単射の
底変更なので、非被約・非平坦な係数変更も許される。
\end{proof}

\begin{remark}[仮定と到達範囲]
\label{ifp:rem:scope}
必要条件の証明で使った因子側の入力は、二次表示と
式\eqref{ifp:eq:two-injections}の二つの単射だけである。
特に因子同士の $j$ 不変量、有限位数・無限位数、二捩れの有無に制約を加えていない。
一方、中心作用を変形する関係や、最高 active 二次関係そのものを動かす変形は
式\eqref{ifp:eq:full-presentation}の対象ではない。
例えば多項式因子では次数 $-1$ の導分が存在し、
$\C\langle x,y,u\rangle/([u,x],[u,y],xy-yx-u^2)$ は
$\C[x]\otimes\C[y]$ を境界に持つ PBW 変形でありながら因子間補正を持つ。
従って固定境界についての低次数単射は実質的な入力である。
本定理は周囲代数の関係を因子別表示へ戻す。
指定中心三次の持上げ、その商の smoothness、可換 Ku 標的の存在は、
それぞれの定理の仮定を別に確認して適用する。
全ての内在的非可換三次がこの固定 tensor 境界を持つとは主張していない。
\end{remark}

\begin{corollary}[一般の二次代数に対する因子分解]
\label{ifp:cor:abstract-factorization}
$V_i$ を有限次元 $\C$ ベクトル空間、
$S_i=T_{\C}(V_i)/(\mathcal R_i)$ を connected な二次代数とし、各因子が
\[
             \operatorname{Der}_{-1}(S_i)=0,
                         \qquad Z(S_i)_1=0
\]
を満たすとする。任意の可換 $\C$-代数 $R$ と任意個の厳密な中心一次変数について、
式\eqref{ifp:eq:full-presentation}、\eqref{ifp:eq:pbw}と同じ意味で
$\bigotimes_iS_i$ の全二次 PBW 変形を考える。
このとき定理\ref{ifp:thm:factorization}の消滅
\eqref{ifp:eq:vanishing}と、元の生成元を保つ tensor 分解
\eqref{ifp:eq:tensor-factorization}が成り立つ。
各因子は PBW 変形であり、結論は任意の係数変更と両立する。
ここでは因子の Koszul 性、AS 正則性、滑らかさを仮定しない。
\end{corollary}

\begin{proof}
式\eqref{ifp:eq:two-injections}で定義した $d_i$ と $\zeta_i$ の核は、
一般の二次代数についても、それぞれ
$\operatorname{Der}_{-1}(S_i)$ と $Z(S_i)_1$ である。
仮定により両写像は有限次元 $\C$ 線型単射なので左逆を持ち、
任意の $R$ への tensor 後にも単射である。
補題\ref{ifp:lem:coefficient-extraction}は PBW 仮定と斉次語基底だけを用いる。
従って式\eqref{ifp:eq:first-commutator}の異なる因子別次数の分離、
内部の他因子出力への $\zeta_i$ の適用、
式\eqref{ifp:eq:second-commutator}の $I^2/I^3$ 係数への $d_i$ の適用が、
一般の $S_i$ についてもそのまま成立する。
因子自身の PBW 性と tensor 表示の同定も、最後の随伴代数の包含による証明を適用できる。
\end{proof}


\begin{corollary}[固定 tensor 境界の全 PBW 中心三次商への比較の適用]
\label{ifp:cor:central-comparison}
定理\ref{ifp:thm:factorization}の表示 $B$ を取り、$k\ge2$、$m\ge1$ とする。
さらに、全て非零の固定重み $a_i\in\C^\times$ に対して、
\[
 w\in Z(B)_3,\qquad
 w\bmod(\mathbf u)=\sum_i a_i\Omega_{i,0}
\]
を満たす指定中心三次を仮定し、$A=B/(w)$ と置く。
このとき次が成り立つ。
\begin{enumerate}
 \item 任意の可換係数環 $R$ 上で、因子別の中心平行移動後、
 $(B,w)$ は命題\ref{emc:prop:coverage}および
 定理\ref{mcp:thm:normal-form}の全許容正規形に入る。
 因子間可換性を変形後の追加仮定として課す必要はない。
 \item $R=\C$ とすると、対応する偏極三次 $X_{\mathrm{pol}}$ に対して
 \[
   \mathcal K_A\text{ が $\C$ 上 smooth}
       \quad\Longleftrightarrow\quad X_{\mathrm{pol}}\text{ が滑らか}
 \]
 である。この条件が成り立てば、次元 $3k+m-2$ の滑らかな可換三次
 $Y_{\mathrm{it}}$ と perfect な実比較核による同値
 \[
                 \Ku_{\mathrm{coh}}(A)\simeq\Ku(Y_{\mathrm{it}})
 \]
 が存在する。
\end{enumerate}
全中心作用を同時に保つ核や、$X_{\mathrm{pol}}$ 自身への比較核を結論するものではない。
\end{corollary}

\begin{proof}
定理\ref{ifp:thm:factorization}により、元の表示で異なる因子は可換であり、
各因子の関係の補正はその因子の生成元と厳密な中心変数だけを含む。
従って一中心では命題\ref{emc:prop:coverage}、多中心では
定理\ref{mcp:thm:normal-form}の入力と正確に一致する。
これらの分類では、二捩れ因子の中心平方方向も含めて
指定境界を持つ全ての中心三次持上げを扱っている。
定理\ref{mcp:thm:normal-form}の小節は $m\ge2$ として記述したが、
その証明の普遍係数代入、平行移動、中心次数ごとの分離は $m=1$ にも同じく適用できる。
従って一中心についても任意の $R$ 上でこの分類を用いてよい。
正規化は因子別の中心平行移動であり、係数環の被約性を要求しない。
これで第一の主張が従う。

複素係数について定理\ref{mcb:thm:smooth-locus}を適用すると、
表示した smoothness の必要十分条件を得る。
その開集合では定理\ref{mcb:thm:central-flags}により一般中央旗が存在し、
系\ref{mcb:cor:all-smooth-comparison}が全 coherent Ku 同値と
perfect 比較核を与える。一中心は定理\ref{euk:thm:universal-kernel}による。
ここで新たに除いた仮定は変形後の因子間可換性であり、
固定 tensor 境界、厳密な中心変数、PBW 性、指定中心三次の存在は保持している。
\end{proof}

% !TEX root = ../../mixed-sklyanin-cubics.tex
% Current result: 2026-11-22 / v61.
\subsection{負次数 Hochschild 類と境界変形の一次中心障害}
\label{ifh:subsec:hochschild}

定理\ref{ifp:thm:factorization}の有限表示に関する結論には、
Hochschild 複体上でも純因子性を説明できる。
ここではその計算を記し、さらに、固定境界の PBW 変形とは別に、
境界自体を次数零で変形するときの一次障害を調べる。
後者は一次の主張であり、形式変形全体や大域的な表示の分類ではない。

各因子を $S_i=T(V_i)/(R_i)=S(a_i',b_i',c_i')$ と書き、
$S=\bigotimes_{i=1}^kS_i$、$V=\bigoplus_iV_i$ と置く。
全因子の標準点曲線は滑らかとし、特に $c_i'\ne0$ である。
中心三次の重み $a_i$ と Sklyanin 係数 $a_i'$ は区別する。
全関係空間は
\[
 \mathsf R=\bigoplus_iR_i\ \oplus\!
   \bigoplus_{i<j}\langle x\otimes y-y\otimes x:
                    x\in V_i,\ y\in V_j\rangle
\]
である。Hochschild 群の下添字は active 内部次数を表す。
以下の複体・群は affine graded 代数 $S$ のものであり、
残余圏の自然変換空間との同一視は用いない。

\begin{lemma}[負次数の純因子性と全 gauge]
\label{ifh:lem:negative}
二捩れ因子を含む任意の固定した因子の組に対して
\begin{equation}\label{ifh:eq:negative-hh}
 \begin{aligned}
 \operatorname{HH}^2(S)_{-1}
   &\cong\bigoplus_i
      \bigl(\operatorname{Sym}_3(\C)/\operatorname{im}T_i\bigr),\\
 \operatorname{HH}^2(S)_{-2}
   &\cong\bigoplus_iR_i^\vee
 \end{aligned}
\end{equation}
となる。ここで $T_i$ は式\eqref{ebc:eq:translation}の平行移動写像であり、
右辺の各直和因子は三次元である。

さらに、quadratic Koszul 双加群複体によりこれらを計算するとき、
次数 $-1$ の二 cocycle 自体は、各 $R_i\to V_i$ の対称行列を並べたものに限る。
次数 $-2$ の二 cocycle 自体も、各 $R_i\to\C$ を並べたものに限る。
いずれも因子間交換関係で零であり、前者は他因子の生成元を値に取らない。

中心変数 $\mathbf u$ を固定し、$\mathbf u=0$ での同定を恒等とする
次数付き $R[\mathbf u]$-代数同型の生成元上の形は、正確に
\begin{equation}\label{ifh:eq:graded-gauge}
 x_{ij}\longmapsto x_{ij}+s_{ij}(\mathbf u),\qquad
              s_{ij}\in R[\mathbf u]_1
\end{equation}
である。上の cochain 計算とこの gauge の記述は、
被約性を仮定しない任意の可換 $\C$-代数 $R$ に両立する。
\end{lemma}

\begin{proof}
一因子の Koszul 双加群解消を
$S_i\otimes K_{i,q}\otimes S_i$ と書く。
$K_{i,q}$ の内部次数は $q$ であり、
$K_{i,0}=\C$、$K_{i,1}=V_i$、$K_{i,2}=R_i$ である。
従って、Hochschild cochain は
\[
 C_i^q=\operatorname{Hom}_{\C}(K_{i,q},S_i),\qquad
 (C_i^q)_d=\operatorname{Hom}_{\C}
               (K_{i,q},(S_i)_{q+d})
\]
で計算できる。特に $\operatorname{HH}^2(S_i)_d=0$ for $d<-2$ である。
この Koszul 複体と PBW 一次条件の関係は
\cite[Sections~3.3, 3.6 and 4.3]{CMOBravermanGaitsgory}に明示されている。

補題\ref{ifp:lem:low-degree}により、
次数 $-1$ の導分と次数一の中心は零である。
次数 $<-1$ の導分も生成元の次数から零であり、従って
\[
       \operatorname{HH}^1(S_i)_{<0}=0,\qquad Z(S_i)_1=0.
\]

各因子の有限自由な Koszul 双加群解消を tensor すると、
$S$ の双加群解消になる。signed shuffle によるその次数二の同定は、
上に表示した $\mathsf R$ の直和分解である。
Hom 複体も因子の Hom 複体の tensor に同定されるため、
体上の複体の K\"unneth 公式から
\begin{equation}\label{ifh:eq:kunneth}
 \operatorname{HH}^*(S)
      \cong\bigotimes_i\operatorname{HH}^*(S_i)
\end{equation}
を cohomological 次数と内部次数の両方を保って得る。
この有限生成自由解消の下での一般公式は
\cite[Theorem~3.1.2]{CMOWitherspoonHochschild}にもある。
固定した次数で現れる項は有限個である。
cohomological 次数二に現れる分割は $(2,0,\ldots)$ と $(1,1,0,\ldots)$ だけである。
前段の消滅と $\operatorname{HH}^0(S_i)=Z(S_i)$ を使うと、内部次数 $-1,-2$ では
純因子以外の項はない。$-1$ の計算は補題\ref{ebc:lem:translation-gauge}である。
$-2$ では $\operatorname{Hom}(R_i,\C)$ の全てが cocycle となる。
唯一の次数三 overlap $\sum x_jr_j=\sum r_jx_j$ に定数を代入すると、
左右が項ごとに相殺するためである。さらに $(C_i^1)_{-2}=0$ なので境界はない。
これで式\eqref{ifh:eq:negative-hh}を得る。

この群の計算だけで実際の関係式を同定するのではなく、境界も確認する。
全体の $(C^1)_{-1}=\operatorname{Hom}(V,\C)$ の coboundary は、
因子間交換関係では零であり、$R_i$ 上では $T_i$ の像となる。
従って全ての coboundary も純因子である。
式\eqref{ifh:eq:kunneth}における純因子類は、
他因子では単位を取ることにより、純因子 cocycle で代表できる。
任意の cocycle とその代表との差も純因子境界なので、
次数 $-1$ の cocycle 自体が主張した支持を持つ。
次数 $-2$ については $(C^1)_{-2}=0$ が同じ結論を与える。
これが有限 PBW 因子分解の cochain 上の理由である。
二次・高次条件を満たす有限表示の結論は、
定理\ref{ifp:thm:factorization}の直接証明による。

最後に、PBW 表示の次数一部分は $V_R\oplus R[\mathbf u]_1$ である。
指定した中心変数を保ち、境界で恒等となる斉次同型は、
従って式\eqref{ifh:eq:graded-gauge}以外の形を持たない。
逆向きの生成元変更は $s_{ij}$ を引く操作である。
これは任意の平行移動が同じ係数の代数の自己同型になるという意味ではなく、
平行移動により変換された係数の表示との同型を記述している。
上で使った各固定次数の複体は有限次元 $\C$-空間の複体なので、
その核・像の計算は任意の $R$ への tensor と両立する。
中心一次式・二次式の各係数にもこの計算を適用できる。
\end{proof}

次に、$\mathbf u=0$ の境界そのものの次数零変形を考える。
この場合には因子間の二次関係を変える混合類が実際に存在する。
まず、その記述に必要な次数零導分を確定する。

\begin{lemma}[次数零導分は Euler 導分だけである]
\label{ifh:lem:euler}
滑らかな標準点曲線を持つ $S_i$ に対し
\[
            \operatorname{HH}^1(S_i)_0=\C E_i
\]
となる。ただし $E_i$ は active 次数による Euler 導分である。
この補題では二捩れを除外しない。
\end{lemma}

\begin{proof}
次数零導分 $D$ は $V_i$ 上の線型写像であり、
その tensor 代数への延長は $R_i$ を保つ。
従って $1+\epsilon D$ は $\C[\epsilon]/(\epsilon^2)$ 上で
関係空間を保つ線型生成元変更である。
標準 Sklyanin 関係の multilinear 行列は
\[
 M(x,y,z)=
 \begin{pmatrix}
 c_i'x&b_i'z&a_i'y\\
 a_i'z&c_i'y&b_i'x\\
 b_i'y&a_i'x&c_i'z
 \end{pmatrix}.
\]
生成元と関係基底の可逆変更は、その determinant を座標変換と単元倍で移す。
従って $D$ の誘導する線型ベクトル場は、
滑らかな点曲線の方程式 $f_i=\det M$ を scalar 倍に保つ。
ある $\lambda\in\C$ に対し $D(f_i)=\lambda f_i$ と書ける。

Euler ベクトル場の $\lambda/3$ 倍を引くと、一次式 $\ell_x,\ell_y,\ell_z$ が
\[
             \ell_x\partial_xf_i+
             \ell_y\partial_yf_i+
             \ell_z\partial_zf_i=0
\]
を満たす。滑らかさと Euler の公式から、三つの二次式 $\partial_jf_i$ の
共通零点は原点だけであり、これらは $\C[x,y,z]$ の正則列である。
その一次 syzygy は Koszul syzygy で生成され、係数の最小次数は二である。
よって上の一次係数は全て零となり、$D$ は Euler 導分の scalar 倍である。
次数零の inner derivation は $(S_i)_0=\C$ から来るので零であり、主張を得る。
\end{proof}

\begin{proposition}[次数零混合類の中心障害とその核]
\label{ifh:prop:central-mixed-obstruction}
二捩れ因子の集合を $I_2$ とし、非零重み $a_i$ と指定中心
\[
                    w_0=\sum_i a_i\Omega_i\in Z(S)_3
\]
を固定する。次の三つの部分空間を置く：
\[
 \begin{aligned}
 \mathcal P&=\bigoplus_i\operatorname{HH}^2(S_i)_0,\qquad
 \mathcal E=\bigoplus_{i<j}\C(E_i\smile E_j),\\
 \mathcal C&=\bigoplus_{\substack{i\ne j\\j\in I_2}}
                  \operatorname{HH}^2(S_i)_{-2}\otimes Z(S_j)_2.
 \end{aligned}
\]
このとき
\begin{equation}\label{ifh:eq:degree-zero}
 \operatorname{HH}^2(S)_0=\mathcal P\oplus\mathcal E\oplus\mathcal C,
           \qquad \dim\mathcal C=9|I_2|(k-1).
\end{equation}

次数付き一次平坦変形 $S_\epsilon$、$\deg\epsilon=0$ に対する
中心持上げの障害写像
\[
 \operatorname{ob}_{w_0}:
      \operatorname{HH}^2(S)_0\longrightarrow\operatorname{HH}^1(S)_3
\]
は
\begin{equation}\label{ifh:eq:mixed-obstruction}
 \operatorname{ob}_{w_0}(E_i\smile E_j)
       =3\bigl(a_i\Omega_iE_j-a_j\Omega_jE_i\bigr)
\end{equation}
を満たし、$\mathcal E$ 上では単射、$\mathcal C$ 上では零である。
$\mathcal E$ の像は純因子変形の障害の像と交わらない。従って
\begin{equation}\label{ifh:eq:obstruction-kernel}
 \ker\operatorname{ob}_{w_0}
   =\left[\bigoplus_i
       \ker\bigl(\operatorname{ob}_{\Omega_i}:
             \operatorname{HH}^2(S_i)_0
                \longrightarrow\operatorname{HH}^1(S_i)_3\bigr)\right]
       \oplus\mathcal C.
\end{equation}
ここで中心補正 $w_0+\epsilon z$ の $z\in S_3$ は因子間を混ぜてよい。
特に全因子が二捩れでない場合、中心を持ち上げられる次数零一次変形の類は
純因子直和 $\mathcal P$ に属する。
\end{proposition}

\begin{proof}
系\ref{eco:cor:non-clifford-lift}と
命題\ref{eco:prop:clifford-center-freedom}の低次数中心の計算により、
各因子で $Z(S_i)_1=0$、$Z(S_i)_3=\C\Omega_i$ である。
さらに $Z(S_i)_2$ は非二捩れ因子では零、二捩れ因子では三つの中心平方が張る。
式\eqref{ifh:eq:kunneth}と補題\ref{ifh:lem:negative}、\ref{ifh:lem:euler}を使う。
cohomological 次数二、内部次数零では、純因子の負次数と
他因子の正次数中心の積は $\operatorname{HH}^2(S_i)_{-2}\otimes Z(S_j)_2$
だけである。これが $\mathcal C$ を与え、各 ordered pair の次元は $3\cdot3=9$ である。
二つの次数一 Hochschild 群を掛ける場合も、負次数がないため、
Euler 導分二つの積だけが残る。これが式\eqref{ifh:eq:degree-zero}である。
同様に
\begin{equation}\label{ifh:eq:target-splitting}
 \begin{aligned}
 \operatorname{HH}^1(S)_3
    ={}&\bigoplus_i\operatorname{HH}^1(S_i)_3
       \ \oplus\!\bigoplus_{i\ne j}\C(\Omega_iE_j)\\
     &\oplus\!\bigoplus_{\substack{i\ne j\\j\in I_2}}
                 \operatorname{HH}^1(S_i)_1\otimes Z(S_j)_2
 \end{aligned}
\end{equation}
を得る。混合項は ordered pair ごとの直和であり、
$\Omega_iE_j$ の内部多重次数 $(3,0)$ と最後の項の $(1,2)$ は異なる。

一次変形の積を $ab+\epsilon\phi(a,b)$ と書くと、結合則は
$\phi$ が Hochschild 二 cocycle であることに等しい。
$w_0$ の中心性から
\[
                  \theta_\phi(b)
                       =\phi(w_0,b)-\phi(b,w_0)
\]
は次数三の導分である。$w_0+\epsilon z$ が中心である条件は
\[
                    \theta_\phi(b)=bz-zb
\]
なので、そのような $z\in S_3$ の存在は $[\theta_\phi]=0$ に等しい。
二 cocycle を coboundary だけ変えたときも、$\theta_\phi$ は inner derivation
だけ変わる。従って、これは変形類の上で定まる障害写像であり、
Gerstenhaber bracket の規約では $[\phi,w_0]$ である。

混合類の代表として
$\phi(a,b)=E_i(a)E_j(b)$ を取る。
$E_i(w_0)=3a_i\Omega_i$ と各 $\Omega_i$ の中心性から、
式\eqref{ifh:eq:mixed-obstruction}を直接得る。
その二項は式\eqref{ifh:eq:target-splitting}の別々の ordered summand に入る。
異なる unordered pair は別の summand を使い、$a_i\ne0$ なので、
従って $\mathcal E$ 上で単射となる。
一方、純因子変形はその因子だけを変形した tensor 積で代表でき、
その中心障害は $\operatorname{HH}^1(S_i)_3$ に入る。
従って両者の相殺はない。

最後に $\beta_i\in\operatorname{HH}^2(S_i)_{-2}$ を取り、
定数補正の単位正規化した一次 cocycle を $\phi_i$ とする。
非二捩れでは系\ref{eco:cor:non-clifford-lift}、二捩れでは
命題\ref{eco:prop:clifford-center-freedom}により、
この定数 PBW 変形には指定中心の持上げがある。
次数二の変形変数を $\delta$ とすると、一次の持上げは
$\Omega_i+\delta q_i$、$q_i\in(S_i)_1$ と書ける。
従って $[\phi_i,\Omega_i]$ は $q_i$ による inner derivation であり、
$\operatorname{HH}^1(S_i)_1$ で零である。

$z_j\in Z(S_j)_2$、$j\ne i$ に対し、
$\phi_i$ を他因子の積で延長し、その値に $z_j$ を掛ける。
これが $\beta_i\otimes z_j\in\mathcal C$ の次数零代表である。
その $w_0$ に対する障害は $a_i z_j[\phi_i,\Omega_i]$ であり、
次数三元 $a_iq_i z_j$ による inner derivation となる。
他因子の $\Omega_l$ はその代表で変形されず、この式には寄与しない。
よって $\mathcal C$ 上の障害は零であり、
式\eqref{ifh:eq:obstruction-kernel}が従う。
\end{proof}

\begin{remark}[固定境界の有限分類と一次障害の区別]
\label{ifh:rem:scope}
負次数の補題は、指定した tensor 境界を保ち、中心変数の可換性を固定した
有限 quadratic PBW 表示に関する定理\ref{ifp:thm:factorization}を説明する。
一方、命題\ref{ifh:prop:central-mixed-obstruction}は、境界自体を変える
次数零一次変形と指定中心三次の持上げに関するものである。
後者から、形式的または大域的な因子分解、全 intrinsic 対象の表示被覆、
あるいは変形の全段階における核の降下を結論しない。
二捩れ因子を許す場合の混合 kernel $\mathcal C$ は明示されたが、
ここで示したのは一次変形に対する中心持上げ障害の消滅である。
その全方向を任意に組み合わせた変形の同時積分可能性や、
全高次の中心持上げまでを含む主張ではない。
\end{remark}

% !TEX root = ../../mixed-sklyanin-cubics.tex
% Current result: 2026-11-22 / v61.
\subsection{Clifford 中心平方による次数零混合類の実積分}
\label{ifc:subsec:clifford-family}

命題\ref{ifh:prop:central-mixed-obstruction}の障害核に現れる
$\operatorname{HH}^2(S_i)_{-2}\otimes Z(S_j)_2$ のうち、
一つの ordered pair に対応する九次元部分は、全係数上の二次代数族へ積分できる。
その構成には、定数型 PBW 因子の普遍係数を、他因子の中心平方へ写す。
中心三次も同じ底変更で持ち上がり、その商は係数環上自由となる。
この族では境界代数自体の次数零変形を行う。
定理\ref{ifp:thm:factorization}で固定した最高 active 次数の tensor 境界も変化している。

\subsubsection{九変数族と中心三次商}

$S_1=S(A_1,B_1,C_1)$ を任意の滑らかな標準 Sklyanin 因子とし、
生成元を $\mathbf x=(x_1,x_2,x_3)$、三つの巡回関係を $r_{1,a}$ と書く。
第二因子を
\[
 S_2=S(1,1,\tau),\qquad
 \tau\ne0,\quad \tau^3\ne1,-8
\]
と固定し、その生成元を $\mathbf y=(y_1,y_2,y_3)$ とする。
このとき
\[
 \{y_2,y_3\}+\tau y_1^2=0,\qquad\text{巡回式},\qquad
 Z(S_2)_2=\langle y_1^2,y_2^2,y_3^2\rangle
\]
である。各因子の指定中心正則三次を $\Omega_{i,0}$ とし、
$a_1,a_2\in\C^\times$ を固定する。
独立な九変数 $\Lambda_{ab}$ の係数環と中心次数二の元を
\begin{equation}\label{ifc:eq:parameters}
 R=\C[\Lambda_{ab}:1\le a,b\le3],\qquad
 z_a=\sum_{b=1}^3\Lambda_{ab}y_b^2\in Z(R\otimes_\C S_2)_2
\end{equation}
と置く。$\Lambda_{ab}$ の内部次数は零である。

\begin{theorem}[中心平方を代入した二次族と自由中心商]
\label{ifc:thm:family}
次の次数付き $R$-代数を考える：
\begin{equation}\label{ifc:eq:presentation}
 \mathcal B_\Lambda
 =R\langle\mathbf x,\mathbf y\rangle
   \big/
    \bigl(r_{1,a}-z_a,\ r_{2,b},\
                   [x_a,y_b]:1\le a,b\le3\bigr).
\end{equation}
全生成元の内部次数を一とする。
この族は $R$ 上自由で、$\Lambda=0$ のファイバーは $S_1\otimes_\C S_2$ である。
各 $S_i$ の固定した斉次語基底の積が $\mathcal B_\Lambda$ の $R$-基底となる。

固定した一次式 $\ell_a(\mathbf x)$ を用いて、中心三次
\begin{equation}\label{ifc:eq:center-lift}
 \Omega_{1,\Lambda}
   =\Omega_{1,0}+\sum_{a,b}\Lambda_{ab}y_b^2\ell_a(\mathbf x)
       \in Z(\mathcal B_\Lambda)_3
\end{equation}
を選べる。ここで $\Omega_{1,0}$ には固定した三次の語代表を用いる。
また $\Omega_{2,0}$ も $\mathcal B_\Lambda$ の中心元である。
従って
\[
 w_\Lambda=a_1\Omega_{1,\Lambda}+a_2\Omega_{2,0},
       \qquad \mathcal A_\Lambda=\mathcal B_\Lambda/(w_\Lambda)
\]
が定まり、$w_\Lambda$ は左右の非零因子である。
$\mathcal A_\Lambda$ も各内部次数で有限自由な $R$-加群であり、Hilbert 級数は
\begin{equation}\label{ifc:eq:hilbert}
 H_{\mathcal B_\Lambda}(t)=\frac1{(1-t)^6},\qquad
 H_{\mathcal A_\Lambda}(t)=\frac{1-t^3}{(1-t)^6}.
\end{equation}
これらの基底・中心性・正則性は、任意の可換 $\C$-代数への係数変更と両立する。
\end{theorem}

\begin{proof}
第一因子の普遍定数型 PBW 代数を
\[
 C_\alpha=\C[\alpha_1,\alpha_2,\alpha_3],\qquad
 \mathcal U_1=C_\alpha\langle\mathbf x\rangle/
                         (r_{1,a}-\alpha_a:1\le a\le3)
\]
とする。$\deg\alpha_a=2$ とし、$\alpha_a$ は中心とする。
一因子の普遍 PBW 基底と中心持上げは、
非二捩れの場合には命題\ref{eco:prop:universal-lift}と
系\ref{eco:cor:non-clifford-lift}、二捩れの場合には
命題\ref{eco:prop:clifford-center-freedom}による。
前者では線形補正を零にして Rees 変数を一へ特殊化すればよい。
従って $\mathcal U_1$ は $C_\alpha$ 上自由で、
$S_1$ の固定した斉次語基底が基底を与える。
また、係数重み二の $\alpha$ と生成元次数一に関する斉次性から、
その中心持上げは
\[
              \Omega_1(\alpha)
                  =\Omega_{1,0}+\sum_a\alpha_a\ell_a(\mathbf x)
\]
と書ける。二捩れの第一因子についても、普遍 Clifford 体積元の同じ重みを用いる。

式\eqref{ifc:eq:parameters}は中心への次数付き準同型
$C_\alpha\to Z(R\otimes S_2)$、$\alpha_a\mapsto z_a$ を与える。
そこで
\begin{equation}\label{ifc:eq:central-base-change}
           \mathcal U_1\otimes_{C_\alpha}(R\otimes_\C S_2)
\end{equation}
を取る。第二因子は非可換だが、$C_\alpha$ は両因子の中心へ写るので
$(p\otimes s)(q\otimes t)=pq\otimes st$ は balanced tensor 積上で定義される。
この代数の普遍性から、式\eqref{ifc:eq:presentation}と同型である。
$\mathcal U_1$ の $C_\alpha$-自由基底を tensor すると主張した基底を得る。
特に $R\otimes S_2$ の $C_\alpha$ 上の平坦性を仮定する必要はない。

$\Omega_1(\alpha)\otimes1$ は第一因子と可換であり、
tensor 積の第二因子とも可換である。その像が
式\eqref{ifc:eq:center-lift}である。
$1\otimes\Omega_{2,0}$ も同様に中心的である。
従って、点ごとに中心元を選び直すことなく、全係数上の中心恒等式を得る。

正則性には、$\mathbf x$ のみを次数一、
$R\otimes S_2$ を次数零とする非負濾過 $F$ を使う。
第一因子の PBW 基底から
\begin{equation}\label{ifc:eq:filtered-algebra}
          \operatorname{gr}_F\mathcal B_\Lambda
                 \simeq S_1\otimes_\C S_2\otimes_\C R
\end{equation}
であり、$w_\Lambda$ の最高 $F$-次数部分は
$a_1\Omega_{1,0}\otimes1$ である。
$\Omega_{1,0}$ は $S_1$ 上の中心正則元であるため、
この最高次部分は式\eqref{ifc:eq:filtered-algebra}上で左右正則である。
有限な最高次数を持つ元の積の初項を比較すれば、
$w_\Lambda$ 自身も左右正則と分かる。

さらに、主イデアルの濾過も strict である。
実際、非零元 $b$ に対し
$\deg_F(w_\Lambda b)=3+\deg_F b$ が成り立ち、
中心性によって二側イデアルの全元は $w_\Lambda b$ の形に書ける。
従って
\begin{equation}\label{ifc:eq:filtered-quotient}
 \operatorname{gr}_F\mathcal A_\Lambda
        \simeq (S_1/(\Omega_{1,0}))\otimes_\C S_2\otimes_\C R.
\end{equation}
$S_1/(\Omega_{1,0})$ の斉次語基底を持ち上げたものと、
$S_2$ の斉次語基底との積を取る。
式\eqref{ifc:eq:filtered-quotient}と最高 $F$-次数に関する帰納法により、
それらは $\mathcal A_\Lambda$ の $R$-自由基底となる。
全ての関係と $w_\Lambda$ は通常の内部次数について斉次なので、
この基底は内部次数別にも斉次である。
各 Sklyanin 因子の Hilbert 級数と $\deg w_\Lambda=3$ から
式\eqref{ifc:eq:hilbert}を得る。
全ての基底表示と初項の単射性は $\C$ 上の自由空間を用いているため、
非被約な係数環への変更にも同じ証明が適用できる。
\end{proof}

\subsubsection{接類と一次障害核への対応}

\begin{proposition}[九つの混合方向の独立性]
\label{ifc:prop:tangent}
$T=S_1\otimes S_2$ とし、二因子次数を
$\deg\mathbf x=(1,0)$、$\deg\mathbf y=(0,1)$ とする。
族 $\mathcal B_\Lambda$ の原点における接写像は同型
\begin{equation}\label{ifc:eq:tangent}
 \operatorname{Mat}_3(\C)
   \xrightarrow{\ \sim\ }
 \operatorname{HH}^2(T)_{(-2,2)}
   \simeq \operatorname{HH}^2(S_1)_{-2}\otimes Z(S_2)_2
       \subset\operatorname{HH}^2(T)_0
\end{equation}
を与える。定数補正 $r_{1,a}-\alpha_a$ に対応する類を $\eta_a$ とすれば、
行列 $M$ の像は $\sum_{a,b}M_{ab}\eta_a\otimes y_b^2$ である。
従って九方向は全て非零かつ独立であり、
命題\ref{ifh:prop:central-mixed-obstruction}の
$\ker\operatorname{ob}_{w_0}$ に属する。
本族はこの一つの ordered pair の九方向を同時に全次数へ積分する。
\end{proposition}

\begin{proof}
式\eqref{ifh:eq:kunneth}の Koszul 双加群複体を二因子次数で分解する。
$R_1$ を第一因子の関係空間とすると、次数の非負性から
\[
 C^1_{(-2,2)}=0,\qquad
 C^2_{(-2,2)}=\operatorname{Hom}_\C(R_1,(S_2)_2)
\]
である。後者の写像を $\psi$ とする。
第一因子だけの三次 overlap では、cocycle 条件は
$\sum_a[x_a,\psi(r_{1,a})]=0$ であり、異なる因子の可換性から自動的に成立する。
二次の第一因子と一次の第二因子の shuffle overlap では、
条件は
\[
                [y,\psi(r)]=0
             \qquad(y\in(S_2)_1,\ r\in R_1)
\]
である。他の三次成分は第一因子の出力次数が負となり寄与しない。
従って cocycle 空間は正確に
$\operatorname{Hom}_\C(R_1,Z(S_2)_2)$ となる。
$C^1_{(-2,2)}=0$ により境界はなく、
三つの関係双対と三つの独立な中心平方から次元は九である。
これは純因子だけでなく、全 Hochschild cochain に対する計算である。

独立パラメータ族を双対数へ $\Lambda_{ab}=\epsilon M_{ab}$、
$\epsilon^2=0$ と制限する。
定理\ref{ifc:thm:family}の自由基底で一次の乗法を取り出すと、
Koszul 二 cochain の $r_{1,a}$ 上の値は
$\sum_bM_{ab}y_b^2$ であり、他の関係上では零である。
これが式\eqref{ifc:eq:tangent}の同型を与える。
ここで接点は九変数 $\Lambda$ の原点であり、$M$ はその接ベクトルである。
最後に、全係数上で存在する $w_\Lambda$ を同じ双対数へ制限すれば、
$w_0=a_1\Omega_{1,0}+a_2\Omega_{2,0}$ の中心持上げを得る。
従って一次中心障害は零であり、実際にはこの族で全次数の持上げが構成されている。
\end{proof}

\begin{corollary}[自由な中心変数の追加]
\label{ifc:cor:central-variables}
任意個の中心次数一変数 $\mathbf u=(u_1,\ldots,u_m)$ を自由に追加し、
$g(\mathbf u)$ を任意の斉次三次式とする。
このとき
\[
 \mathcal B_\Lambda[\mathbf u],\qquad
 \mathcal A_{\Lambda,g}
   =\mathcal B_\Lambda[\mathbf u]/
       (w_\Lambda+g(\mathbf u))
\]
についても、表示した中心三次は左右正則であり、
$\mathcal A_{\Lambda,g}$ は $R[\mathbf u]$ 上自由である。
各複素係数点での内部 Hilbert 級数は
\[
              H_{\mathcal A_{\Lambda,g}}(t)
                       =\frac{1-t^3}{(1-t)^{6+m}}
\]
である。$g$ の係数も独立な変数として加えてよい。
\end{corollary}

\begin{proof}
前証明の $F$ に関して $\mathbf u$ の次数を零とする。
中心三次の最高 $F$-次数部分は引き続き $a_1\Omega_{1,0}$ である。
式\eqref{ifc:eq:filtered-quotient}の右辺に
$\C[\mathbf u]$ と $g$ の係数環を tensor し、同じ基底の議論を適用すればよい。
\end{proof}

\begin{remark}[この積分族の範囲]
\label{ifc:rem:scope}
本小節で得たものは、固定 Clifford 因子の中心平方を使う、
一方向の九変数混合族とその中心正則三次商である。
式\eqref{ifh:eq:obstruction-kernel}の他の ordered pair、
逆向きの混合類、純因子の次数零変形と同時に積分できるとは結論していない。
また、この $F$ は第一因子だけの次数濾過であり、
固定 tensor 境界の中心一次・二次補正を扱う前小節の active 濾過とは異なる。
PBW 基底、係数環上の自由性、中心持上げとその正則性がここでの結論である。
周囲代数の AS-regular 性・CY 性、射影中心商や残余圏の smoothness、
および可換三次への Ku 比較核は本小節では確認していない。
中心変数を加えない Hilbert 型は四次元三次に対応するため、
主稿の $n\ne4$ という仮定を持つ開性定理
（定理\paperref{main}{grl:thm:open-locus}）をそのまま用いることもできない。
一つの中心変数と非零三乗項を加えた五次元型では、この次元上の制約はなくなるが、
同定理に必要な実 smooth proper 圏族などの入力は別途確認しなければならない。
\end{remark}

% !TEX root = ../../mixed-sklyanin-cubics.tex
% Current result: 2026-11-22 / v62.
\subsection{中心平方混合族の相対 Koszul 性と非捩れ双対性}
\label{ifr:subsec:regularity}

定理\ref{ifc:thm:family}の周囲代数には、全係数で有限な自己双対
Koszul 双加群解消が存在する。ここでは係数環の正則性も被約性も仮定しない。
一方、複素ファイバーの AS 正則性と、係数環上の相対的な双対性は区別する。
Koszul 双対の Frobenius 形式から双加群の Nakayama twist を計算する方法は
\cite[Theorem~6.3, Proposition~6.4]{IFRBergerMarconnet}および
\cite[Lemmas~1.1--1.2]{IFRHeVanOystaeyenZhang}に基づく。
以下では、一般の可換係数環へその結論だけを移すのではなく、
有限自由解消と双対写像を構成する。

\begin{theorem}[全係数での相対正則性と非捩れ Calabi--Yau 双対性]
\label{ifr:thm:ambient-regularity}
定理\ref{ifc:thm:family}の二因子 $S_1,S_2$ を固定する。
$R$ を任意の可換 $\C$-代数、$\Lambda_{ab}\in R$ とし、
同定理の底変更を $B=\mathcal B_\Lambda\otimes_{\C[\Lambda]}R$ と書く。
$B$ の augmentation は全ての $x_a,y_b$ を零に写すものとし、
$B_R^{\mathrm e}=B\otimes_R B^{\mathrm{op}}$ と置く。
次が成り立ち、表示する解消と双対同型は任意の可換係数変更と両立する。
\begin{enumerate}
 \item $B$ は $R$ 上相対 Koszul である。すなわち $V=B_1=R^6$ と
 $\mathcal R_B=\ker(V^{\otimes_R2}\to B_2)$ に対し、
 \[
 K_0=R,\quad K_1=V,\quad
 K_i=\bigcap_{a+b=i-2}
       V^{\otimes_R a}\otimes_R\mathcal R_B\otimes_RV^{\otimes_R b}
           \ \subset V^{\otimes_R i}\quad(i\ge2)
 \]
 は階数 $\binom6i$ の有限自由 $R$-加群であり、$i>6$ では零となる。
 $B\otimes_RK_i$ は $R$ の左線形自由解消を、
 $K_i\otimes_RB$ は右線形自由解消を与える。
 \item $P_i=B\otimes_RK_i\otimes_RB$ を項とする長さ六の
 Koszul 双加群複体 $P_\bullet\to B$ は有限自由
 $B_R^{\mathrm e}$-解消である。従って $B$ は $R$ 上 homologically smooth である。
 \item 内部次数付き双対について
 \begin{equation}\label{ifr:eq:bimodule-duality}
  \mathbf R\underline{\operatorname{Hom}}_{B_R^{\mathrm e}}
              (B,B_R^{\mathrm e})\simeq B(6)[-6]
 \end{equation}
 が成り立つ。右 $B_R^{\mathrm e}$-加群は外側二作用を交換して
 $B$-双加群と見る。右辺に Nakayama twist はない。また両側で
 \begin{equation}\label{ifr:eq:augmentation-duality}
  \mathbf R\underline{\operatorname{Hom}}_B(R,B)\simeq R(6)[-6],
  \qquad
  \mathbf R\underline{\operatorname{Hom}}_{B^{\mathrm{op}}}(R,B)
                 \simeq R(6)[-6]
 \end{equation}
 である。
 \item $R=\C$ の各ファイバーは大域次元六、AS index 六の
 Koszul AS-regular 代数であり、非捩れ $6$-Calabi--Yau 代数である。
 Hilbert 級数は $(1-t)^{-6}$、GK 次元は六である。
\end{enumerate}
\end{theorem}

\begin{proof}
\textbf{第一因子の濾過と、底変更に安定な Koszul 空間。}
$x_a$ の濾過次数を一、$y_b$ の濾過次数を零とする増大濾過を $F$ とする。
定理\ref{ifc:thm:family}の積基底の証明から
\begin{equation}\label{ifr:eq:associated}
 \operatorname{gr}_F B=D\otimes_\C R,
 \qquad D=S_1\otimes_\C S_2
\end{equation}
であり、各内部次数 $d$ では $0=F_{-1}B_d\subset F_0B_d
\subset\cdots\subset F_dB_d=B_d$ という有限濾過になる。
各 associated graded は固定した有限次元 $\C$-空間の $R$ への拡大である。
両 $S_i$ は quadratic Koszul かつ非捩れ $3$-Calabi--Yau であり、
それらの解消の tensor 積から $D$ は Koszul である。
滑らかな標準 Sklyanin 因子の AS 正則性と Koszul 性は既に用いた一因子入力であり、
その巡回 potential 表示と\cite[Corollary~2.9]{MoriSmith}が非捩れ $3$-CY 性を与える。
Clifford 因子も補題\ref{eco:lem:clifford-input}の仮定に入る。

被約 bar 複体 $C_\bullet(B)=\operatorname{Bar}_R(R,B,R)$ を取る。
内部次数 $d$ の部分は
\[
 C_i(B)_d=
  \bigoplus_{d_1+\cdots+d_i=d,\ d_j>0}
          B_{d_1}\otimes_R\cdots\otimes_RB_{d_i}
\]
で、$0\le i\le d$ の有限自由 $R$-加群の複体である。
総 $x$-次数で濾過すると、その associated graded は
$C_\bullet(D)_d\otimes_\C R$ となる。
$D$ の Koszul 性から、その homology は homological degree $d$ だけにあり、
\[
 H_d(C_\bullet(D)_d)=K_d(D),\qquad
 \dim_\C K_d(D)=\binom6d
\]
である。特に $d>6$ では全体が acyclic である。

この有限濾過の議論は非被約環上でも分裂する。
実際、各 $d$ と各 $x$-次数で、$\C$ 上の bar 複体を
その homology と可縮な二項複体の直和へ分け、縮約を一つ固定する。
積基底を用いて $C_\bullet(B)_d$ の底加群を
$C_\bullet(D)_d\otimes R$ と同定すると、微分の差 $\delta$ は
$x$-濾過を真に下げる。固定した縮約の homotopy は $x$-次数を保つので、
摂動縮約に現れる $h\delta,\delta h$ の冪級数は有限和である。
転送された微分は唯一の homological degree $d$ から次数を一下げるため零となる。
従って $C_\bullet(B)_d$ の homology は
$K_d(D)\otimes R$ と濾過付きで同型であり、その associated graded の同型は
式\eqref{ifr:eq:associated}から来る自然なものに一致する。
この構成は $\C$ 上で固定した線形分裂と係数の有限演算だけを用いるので、
任意の $R\to R'$ と両立する。

bar の最上項は $C_d(B)_d=V^{\otimes_Rd}$ で、そこへ入る微分はない。
その微分の核は、隣接二文字の積の全てが零となる部分、すなわち上記の $K_d$ である。
従って各 $K_d$ は $V^{\otimes_Rd}$ の直和因子であり、
\begin{equation}\label{ifr:eq:k-associated}
  \operatorname{gr}_F K_d=K_d(D)\otimes_\C R.
\end{equation}
その各 $x$-重み成分は
$K_p(S_1)\otimes_\C K_{d-p}(S_2)\otimes_\C R$ と同型である。
交叉の形成が一般に底変更と可換だと仮定したのではなく、
以上の分裂した bar 計算によって、ここでの $K_d$ に対してそれを得た。

Koszul 複体の微分を具体的に記す。
$K_i$ の切断余積を
\[
 \Delta_{1,i-1}(k)=\sum_s v_s\otimes k_s,\qquad
 \Delta_{i-1,1}(k)=\sum_s k'_s\otimes v_s
\]
と書けば、双加群微分は
\begin{equation}\label{ifr:eq:koszul-differential}
 d_i(b\otimes k\otimes c)=
  \sum_s bv_s\otimes k_s\otimes c
       +(-1)^i\sum_s b\otimes k'_s\otimes v_sc.
\end{equation}
余積が $K$ の tensor 積に入ることは、上の交叉表示と直和因子性から従う。
二次関係により $d^2=0$ である。
その associated graded は $D$ の Koszul 双加群解消を $R$ へ拡大したものとなる。
各内部次数で有限濾過を用いれば完全性が従う。
同じ議論を片側の複体に適用する。
係数変更後も式\eqref{ifr:eq:associated}と\eqref{ifr:eq:k-associated}が成立するため、
これらは任意の底変更後にも解消である。

\textbf{反転双対濾過と Frobenius 形式。}
$E_i=\operatorname{Hom}_R(K_i,R)$、$E=\bigoplus_{i=0}^6E_i$ と置き、
切断余積の双対を積とする。この相対 Koszul 双対は
通常の次数に関する符号同型を通じて
$\operatorname{Ext}_B^*(R,R)$ に一致する。
包含 $K_i\subset V^{\otimes_Ri}$ は分裂しているので、
$(V^*)^{\otimes_R i}\to E_i$ は全射であり、$E$ は $E_1$ で生成される。
ここでの $E_i$ の次数 $i$ は Ext 次数を表し、内部双対次数は $-i$ である。

$E_i$ には増大する反転双対濾過
\begin{equation}\label{ifr:eq:dual-filtration}
 F^!_qE_i=\{e\in E_i:e(F_{i-q-1}K_i)=0\}
\end{equation}
を入れる。$x_a^*=\xi_a$ の重みは零、$y_b^*=\eta_b$ の重みは一となる。
切断余積の濾過互換性から $F^!_qE_i\,F^!_rE_j\subset F^!_{q+r}E_{i+j}$ であり、
\begin{equation}\label{ifr:eq:dual-associated}
 \operatorname{gr}_{F^!}E
   =(S_1^!\widehat\otimes_\C S_2^!)\otimes_\C R.
\end{equation}
$\widehat\otimes$ は Ext 次数の Koszul 符号を持つ tensor 積を表す。
式\eqref{ifr:eq:k-associated}を双対化すると、より具体的に
\[
 \operatorname{gr}_{q}^{F^!}E_i
  =(S_1^!)_{i-q}\otimes_\C(S_2^!)_q\otimes_\C R
\]
である。従って全ての濾過商は固定された $\C$-空間の拡大であり、有限自由である。

各 $S_j^!$ の次数三の向きを固定する。
$E_6$ には重み三だけが存在するので
\[
 F^!_2E_6=0,\qquad F^!_3E_6=E_6,
 \qquad E_6\xrightarrow{\sim}\operatorname{gr}^{F^!}_3E_6\xrightarrow{\sim}R
\]
により trace $t:E_6\to R$ が自然に定まる。
他次数では $t$ を零とする。この trace は上の向きを固定すれば係数変更と両立する。
$S_j$ の非捩れ $3$-CY 双対性により $S_j^!$ は graded symmetric Frobenius である。
従って式\eqref{ifr:eq:dual-associated}の積対合は、次数 $i,6-i$、
重み $q,3-q$ の間で完全である。

実際の対合
\begin{equation}\label{ifr:eq:frobenius-pairing}
       E_i\otimes_R E_{6-i}\longrightarrow R,\qquad (e,f)\longmapsto t(ef)
\end{equation}
も完全である。これを係数環の閉点でだけ検査する必要はない。
相手側の双対へ
$\operatorname{Ann}(F^!_{2-q}E_{6-i})$ という濾過を入れると、
左および右の対合写像は濾過を保ち、その associated graded は上の完全対合となる。
有限濾過の帰納により両写像は濾過付き同型である。
特に左右いずれの対合についても
\begin{equation}\label{ifr:eq:annihilator}
 F^!_qE_i
  =\{e\in E_i:t(eF^!_{2-q}E_{6-i})=0\}
  =\{e\in E_i:t(F^!_{2-q}E_{6-i}e)=0\}
\end{equation}
となる。これは任意の可換 $R$ 上の有限自由加群の主張である。

\textbf{符号対合による Nakayama shear の消去。}
完全対合により一意に定まる graded Frobenius Nakayama 自己同型 $\mu$ を
\[
                       t(ef)=t(f\mu(e))
\]
で正規化する。結合則と完全性から $\mu$ は代数自己同型であり、
式\eqref{ifr:eq:annihilator}から $F^!$ を保つ。
associated graded の trace は次数六の graded symmetric trace なので、
その Nakayama 自己同型は次数一で $-1$ である。
従って或る $c_{ba}\in R$ に対し
\begin{equation}\label{ifr:eq:possible-shear}
 \mu(\xi_a)=-\xi_a,\qquad
 \mu(\eta_b)=-\eta_b+\sum_a c_{ba}\xi_a
\end{equation}
としかなり得ない。

表示\eqref{ifc:eq:presentation}には実際の自己同型
\[
        \sigma(x_a)=-x_a,\qquad \sigma(y_b)=y_b
\]
がある。各 $r_{1,a}$ と $z_a$ が二次であることがここで必要である。
双対上では $\sigma(\xi_a)=-\xi_a$、$\sigma(\eta_b)=\eta_b$ となり、
唯一の top 重み成分 $(S_1^!)_3\otimes(S_2^!)_3$ 上では $-1$ を掛ける。
従って $t\sigma=-t$ である。
trace を単元倍する自己同型は、その Frobenius Nakayama 自己同型と可換する。
ここでは直接
\[
 t(\sigma(e)f)
    =-t(e\sigma^{-1}(f))
    =-t(\sigma^{-1}(f)\mu(e))
    =t(f\sigma\mu(e))
\]
を用い、完全性から $\mu\sigma=\sigma\mu$ と分かる。
式\eqref{ifr:eq:possible-shear}に代入すると $2c_{ba}=0$ であり、
$R$ は $\C$-代数なので全ての $c_{ba}$ は零となる。
$E$ は $E_1$ で生成されるから、$\mu|_{E_i}=(-1)^i$ である。
従って
\begin{equation}\label{ifr:eq:graded-symmetry}
          t(ef)=(-1)^{ij}t(fe)\qquad(e\in E_i,\ f\in E_j).
\end{equation}
$i+j\ne6$ のとき両辺は零であり、$i+j=6$ では
$i\equiv i(6-i)\pmod2$ を使った。

\textbf{有限自由複体の自己双対。}
最後に、非被約係数環上でも twist がないことを複体上で確かめる。
trace 対合から同型
\[
 \theta_i:E_i\longrightarrow K_{6-i},\qquad
          \theta_i(e)(f)=t(ef)\quad(f\in E_{6-i})
\]
を得る。内部次数の差は六である。
$v_s$ を $V$ の基底、$e_s$ を $E_1$ の双対基底とする。
双加群複体 $P$ の Hom 双対を外側二因子を交換して
$B\otimes_RE_i\otimes_RB$ と書く。
標準 Hom 微分 $\delta_i(f)=(-1)^{i+1}f d_{i+1}$ はこの表示で
\[
 \delta_i(b\otimes e\otimes c)=
  \sum_s bv_s\otimes ee_s\otimes c
       +(-1)^{i+1}\sum_s b\otimes e_se\otimes v_sc
\]
である。一方、$\theta_i(e)$ の左右切断余積は
\[
 \Delta_{1,5-i}\theta_i(e)
       =\sum_s v_s\otimes\theta_{i+1}(ee_s),\qquad
 \Delta_{5-i,1}\theta_i(e)
       =-\sum_s\theta_{i+1}(e_se)\otimes v_s.
\]
ここでは $0\le i\le5$ とする。
第二式では $t(efe_s)=-t(e_sef)$、すなわち
式\eqref{ifr:eq:graded-symmetry}の総次数五と一の交換を用いた。
これを式\eqref{ifr:eq:koszul-differential}へ代入すると
$1\otimes\theta_i\otimes1$ は
\[
 \underline{\operatorname{Hom}}_{B_R^{\mathrm e}}(P_\bullet,B_R^{\mathrm e})
                     \xrightarrow{\sim}P_\bullet(6)[-6]
\]
という複体同型を与える。
従って式\eqref{ifr:eq:bimodule-duality}が従う。
構成に現れる全ての項は有限自由であり、対合と trace は底変更と両立する。

$P_\bullet\to B$ は左または右の $B$-加群として分裂した解消である。
実際、各項はその側で自由で、終項 $B$ も自由なので、短完全列を順に分裂できる。
上の双対同型の外側一因子へ augmentation を施すと、
左 Koszul 解消の Hom 双対が右 Koszul 解消の六次 shift に同定される。
これが式\eqref{ifr:eq:augmentation-duality}を与える。

$R=\C$ では、任意の左 $B$-加群 $M$ について
$P_\bullet\otimes_BM\to M$ は長さ六の自由解消である。
従って大域次元は高々六で、augmentation の第六 Ext が非零なので丁度六となる。
両側の AS-Gorenstein 条件と index は
式\eqref{ifr:eq:augmentation-duality}による。
Hilbert 級数から GK 次元も六であり、AS 正則性が従う。
双加群解消と式\eqref{ifr:eq:bimodule-duality}は非捩れ CY 条件そのものである。
\end{proof}

\begin{corollary}[Noetherian 性と enveloping algebra の相対有限解消]
\label{ifr:cor:noetherian-envelope}
定理\ref{ifr:thm:ambient-regularity}の記号を用いる。
$R$ が有限生成可換 $\C$-代数の局所化なら、$B$ と $B_R^{\mathrm e}$ は左右 Noetherian である。
任意の可換 $R$ について、$B_R^{\mathrm e}$ 自身にも長さ十二の
有限自由な相対 Koszul 双加群解消が存在する。
従って $R$ が大域次元 $r<\infty$ を持つなら
\[
 \operatorname{gldim} B\le6+r,\qquad
 \operatorname{gldim} B_R^{\mathrm e}\le12+r
\]
が左右で成り立つ。
\end{corollary}

\begin{proof}
補題\paperref{main}{mbnp:lem:noetherian}を二因子、および
$S_1,S_2,S_1^{\mathrm{op}},S_2^{\mathrm{op}}$ の四因子に適用する。
同補題の multi-homogeneous coordinate ring と中心正則列による証明は、
Clifford 因子と opposite 因子にも適用できる。
従って $D$ と $D\otimes_\C D^{\mathrm{op}}$ は左右 Noetherian である。
有限生成可換係数環への拡大とその局所化は、中心多項式拡大・商・中心局所化であるから、
\[
 D\otimes_\C R,\qquad
 (D\otimes_\C D^{\mathrm{op}})\otimes_\C R
\]
も左右 Noetherian である。
式\eqref{ifr:eq:associated}と、二つの $x$-濾過の和を用いた
\[
 \operatorname{gr}_F B_R^{\mathrm e}
              =(D\otimes_\C D^{\mathrm{op}})\otimes_\C R
\]
から、先頭項の持上げにより $B,B_R^{\mathrm e}$ の主張が従う。
ここで tensor の associated graded は、各濾過商が $R$ 自由であるため上の表示になる。

$B$ と $B^{\mathrm{op}}$ の有限自由相対 Koszul 双加群解消を $R$ 上で
tensor すれば、$B_R^{\mathrm e}$ の長さ十二の解消を得る。
完全性と任意の底変更は、両解消が $R$-加群として分裂していることから従う。
最後に、$P_\bullet\otimes_BM$ の各項は $B\otimes_RM$ の有限直和である。
$R$-加群 $M$ の長さ $r$ の射影解消を $R$-自由な $B$ へ誘導すれば、
この各項の $B$-射影次元は高々 $r$ となる。
従って長さ六の複体全体から上界 $6+r$ を得る。
enveloping algebra にも長さ十二の解消を使って同じ議論を適用する。
\end{proof}

\begin{corollary}[中心多項式拡大と正則三次商の augmentation 双対性]
\label{ifr:cor:central-extension}
任意の可換 $R$ 上で $m\ge0$ とし、
$B^{(m)}=B[u_1,\ldots,u_m]$、$\deg u_j=1$ と置く。
この代数には任意の底変更と両立する長さ $N=6+m$ の有限自由 Koszul 双加群解消があり、
\[
 \mathbf R\underline{\operatorname{Hom}}_{(B^{(m)})_R^{\mathrm e}}
       (B^{(m)},(B^{(m)})_R^{\mathrm e})\simeq B^{(m)}(N)[-N],
 \qquad
 \mathbf R\underline{\operatorname{Hom}}_{B^{(m)}}(R,B^{(m)})
                      \simeq R(N)[-N]
\]
である。augmentation の式は右側にも成り立つ。
任意の $g\in R[u_1,\ldots,u_m]_3$ に対し、
系\ref{ifc:cor:central-variables}の中心正則元
$\widetilde w=w_\Lambda+g$ と商 $A=B^{(m)}/(\widetilde w)$ を取ると、
\begin{equation}\label{ifr:eq:quotient-augmentation}
 \mathbf R\underline{\operatorname{Hom}}_A(R,A)
       \simeq R(3+m)[-(5+m)],\qquad
 \mathbf R\underline{\operatorname{Hom}}_{A^{\mathrm{op}}}(R,A)
       \simeq R(3+m)[-(5+m)]
\end{equation}
であり、任意の底変更と両立する。
$R$ が有限生成 $\C$-代数の局所化なら
$B^{(m)},(B^{(m)})_R^{\mathrm e},A,A_R^{\mathrm e}$ は左右 Noetherian である。
\end{corollary}

\begin{proof}
定理\ref{ifr:thm:ambient-regularity}の解消と通常の多項式環の Koszul 解消を
$R$ 上で tensor する。top trace の符号付き tensor 積は graded symmetric であり、
その自己双対は表示した $N$ 次元の非捩れ双対性を与える。
$\widetilde w$ の左右正則性と商の $R$-自由性は
系\ref{ifc:cor:central-variables}による。
二項解消 $B^{(m)}(-3)\xrightarrow{\widetilde w}B^{(m)}$ から
\[
 \mathbf R\underline{\operatorname{Hom}}_{B^{(m)}}(A,B^{(m)})
                    \simeq A(3)[-1]
\]
を得る。change-of-rings 随伴により
\[
 \mathbf R\underline{\operatorname{Hom}}_{B^{(m)}}(R,B^{(m)})
   \simeq
 \mathbf R\underline{\operatorname{Hom}}_A(R,A)(3)[-1],
\]
従って式\eqref{ifr:eq:quotient-augmentation}となる。
任意の $R\to R'$ に対して change-of-rings 随伴の底変更図式を取る。
$B^{(m)}$ 側の augmentation 双対の底変更射は有限自由 Koszul 解消により同型であり、
$\mathbf R\underline{\operatorname{Hom}}_{B^{(m)}}(A,B^{(m)})$ の底変更射も
任意の底変更で完全な上記の二項解消により同型である。
従って同じ図式の $A$ 側の augmentation 双対の底変更射も同型となる。
Noetherian 性は系\ref{ifr:cor:noetherian-envelope}から中心多項式拡大と商で従う。
\end{proof}

\begin{remark}[相対結論と射影中心商の滑らかさの区別]
\label{ifr:rem:scope}
定理\ref{ifr:thm:ambient-regularity}は、全 $\Lambda$ に対する周囲代数の結果である。
任意の係数環 $R$ 上での相対 smoothness は、特異な $R$ を基礎とした
$\C$ 上の絶対 smoothness や有限大域次元を主張するものではない。
また、三次商 $A$ の augmentation 双対性だけから
$\mathrm D^b(\operatorname{qgr}A)$ または残余圏の smoothness は従わない。
中心商の相対 perfect 対角、可換 Ku への比較、全混合障害核の同時積分は、
それぞれ別の入力または構成を要する。
\end{remark}

% !TEX root = ../../mixed-sklyanin-cubics.tex
% Current result: 2026-11-22 / v62.
\subsection{Clifford 混合族の実相対対角と滑らかな近傍}
\label{iff:subsec:smoothness}

定理\ref{ifc:thm:family}の九変数族について、分離した原点を含む開集合上で
実際の smooth proper な残余圏族を構成する。
ここでは通常の二重次数 coherent 対角を使う。
特殊ファイバーの有限複体を任意に持ち上げるのではなく、既に全係数上にある
対角の線束解消を有限段で分裂させることが要点である。
四次元型にもこの議論は適用できる。
この小節では $n\ne4$ の可換実現の開性定理を五次元型へ適用し、
四次元型は後続の系\ref{ifk:cor:clifford-fourfold-neighborhood}で
rotation を用いて扱う。

係数環を $R_0=\C[\Lambda_{ab}:1\leq a,b\leq3]$ と書き、
$o=V(\Lambda_{ab})$ とする。
以下の二つの場合を同時に扱う：
\begin{equation}\label{iff:eq:families}
 \begin{array}{c|c|c|c}
  m&B&A&n\\ \hline
  0&\mathcal B_\Lambda&\mathcal B_\Lambda/(w_\Lambda)&4\\
  1&\mathcal B_\Lambda[u]&
      \mathcal B_\Lambda[u]/(w_\Lambda+cu^3)&5
 \end{array}
 \qquad c\in\C^\times .
\end{equation}
$m=1$ の $c$ は固定し、$N=6+m$、$a=N-3$、$n=N-2$ と置く。
第一因子は引き続き任意の滑らかな標準 Sklyanin 因子でよい。
以下で使う代数的入力は、定理\ref{ifr:thm:ambient-regularity}、
系\ref{ifr:cor:central-extension}、\ref{ifr:cor:noetherian-envelope}の
相対 Koszul 解消、augmentation 双対、包絡環の noetherian 性と、
定理\ref{ifc:thm:family}の次数別自由性である。

\subsubsection{proper な線束模型と通常の射影積}

\begin{proposition}[有限線束模型と係数変更]
\label{iff:prop:proper-model}
$E=\bigoplus_{i=0}^{N-1}\mathcal O(i)$ とし、
$G=L_{\mathcal O}\cdots L_{\mathcal O(a-1)}E$ と置く。
全線束の厚生成圏とその残余圏は、それぞれ
\[
 \operatorname{Perf}(\Sigma_{R_0}),\qquad
 \operatorname{Perf}(\Gamma_{R_0}),\qquad
 \Sigma_{R_0}=R\operatorname{End}(E),\quad
 \Gamma_{R_0}=R\operatorname{End}(G)
\]
と表示され、両 dg 代数は $R_0$-perfect である。
有限型可換 $R_0$-代数とその局所化への係数変更は、
これらの模型、例外列、mutation、および内部 shift を対応する模型へ移す。
特に非被約な局所 Artin 環への変更も含む。
底上の導来 tensor による圏族は任意の導来底変更に対して定義される。
この段階では残余線束圏と全 coherent 残余圏を同一視しない。
\end{proposition}

\begin{proof}
上に挙げた入力と中心三次の二項解消から、両側で
\[
 R\underline{\operatorname{Hom}}_B(R_0,B)=R_0(N)[-N],
 \qquad
 R\underline{\operatorname{Hom}}_A(R_0,A)=R_0(a)[-N+1]
\]
を得る。補題\paperref{main}{lp:lem:local-cohomology}の
augmentation 解消による計算は特殊ファイバーの可換性を使わないので、
本族にも適用できる。従って
\begin{equation}\label{iff:eq:line-hom}
 \operatorname{Ext}^q(\mathcal O(i),\mathcal O(j))=
 \begin{cases}
  A_{j-i},&q=0,\\
  (A_{i-j-a})^\vee,&q=n,\\
  0,&\text{その他}
 \end{cases}
\end{equation}
であり、各項は有限自由な係数加群である。
ここで $(-)^\vee$ は係数環上の双対である。
ambient の長さ $N$ の augmentation 解消を $A$ へ tensor すると、
cohomology は $R_0$ と $R_0(-3)$ のみである。
両者は torsion なので tails における有限完全複体を得る。
その両端の線束を順に解けば、連続した $N$ 本の線束が全 twist を厚生成する。
式\eqref{iff:eq:line-hom}と有限 mutation から properness が従う。

係数変更は Hom の次元だけによる判定ではなく、計算複体に対して行う。
ambient の相対解消は内部次数別に係数環上分裂している。
$A$ の二項 ambient 解消も、その余核が次数別自由なので同じ性質を持つ。
augmentation 双対から得られる ambient の導来 saturation と、その二項複体への
適用は、従って通常の任意の係数変更と両立する。
有限個の線束 Hom、evaluation、cone で作る模型と mutation にもこれを適用できる。
これにより Artin 底変更を含む主張を得る。
\end{proof}

$R$ を $R_0$ の正則な局所化、原点の正則局所環を
正則パラメータの一部で割った環、または任意の複素ファイバーの係数環とする。
対応する代数を再び $A,B$ と書き、
\[
 C=A^{\mathrm{op}}\otimes_R A,\qquad
 D=B^{\mathrm{op}}\otimes_R B
\]
に二重内部次数を入れる。
$\operatorname{QGr}^{(2)}C$ は、各元が十分大きい
$C_{\geq r,\geq r}$ によって消える加群を除いた商とする。
その noetherian 対象を $\operatorname{qgr}^{(2)}C$ と書く。
これは \cite[Definitions~3.16, 3.18]{CMOBallardFarman} の通常の射影積である。

\begin{lemma}[正則な係数環上の射影積の有限性]
\label{iff:lem:projective-product}
この射影積では、線束 $\mathcal O(i,j)$ が非有界導来圏の
compact generator の集合をなし、任意の通常の対象 $Q$ に対して
\[
 \operatorname{Ext}^q(\mathcal O(i,j),Q)=0
       \qquad(q>2(N-1))
\]
である。coherent な $Q$ に対する線束からの導来 Hom は $R$-perfect であり、
coherent な $M,Q$ に対する各 $\operatorname{Ext}^q(M,Q)$ は有限 $R$-加群である。
これらの Ext は中心的な局所化と両立する。
$Q$ が次数別 $R$-flat な graded 加群で表示される場合、線束 Hom の計算は
正則パラメータの商を含む係数変更と両立する。
一変数の射影圏についても同じ結論が成り立ち、上界は $N-1$ でよい。
\end{lemma}

\begin{proof}
noetherian 性は $B$ 自身だけでなく $D$ にも必要であり、
系\ref{ifr:cor:noetherian-envelope}を使う。
その根拠は主稿の補題\paperref{main}{mbnp:lem:noetherian}の
多重 THCR の証明を、二つの因子とその opposite の四因子へ適用することである。
個々の因子の strong noetherian 性だけから tensor 積の noetherian 性を推論しない。

$B_{<r}=B/B_{\geq r}$ と置く。
中心三次商の二項解消による
$\operatorname{Tor}_1^B(A,B_{<r})$ は内部次数 $r,r+1,r+2$ のみにあり、
$r'\geq r+3$ の遷移は零である。
従って graded $A$-injective で計算した torsion の正次数部分は
直極限で消え、$A$-加群上で
\begin{equation}\label{iff:eq:torsion-independence}
 R\tau_A=R\tau_B,\qquad RQ_A=RQ_B
\end{equation}
となる。これは係数環上の flatness と中心元の正則性による計算である。
相対対角解消を $R$-free な $B_{<r}$ へ tensor すると、
$\operatorname{pd}_B B_{<r}\leq N$ を得る。
noetherian 性により有限生成射影項で計算できるから、
$R\tau_B$ は coproduct と可換であり、cohomological dimension は高々 $N$ である。
$RQ_B$ の正次数は $N-1$ 以下となる。

二つの truncation による saturation の合成は
\cite[Proposition~3.41]{CMOBallardFarman} の計算と同じく
\[
 RQ_C=RQ_{A^{\mathrm{op}}}\,RQ_A
       =RQ_{B^{\mathrm{op}}}\,RQ_B
\]
を与える。相対計算では tensor を導来 tensor とする。
両因子の $R$-flatness は restriction が injective を保つために使い、
truncation の有限射影解消と coproduct の可換性は直前の議論を使う。
この合成の内部次数成分が線束からの Hom なので、表示した上界と compact 性を得る。
全成分が零なら導来 saturation 自体が零であるため、生成性も成り立つ。

$d=\operatorname{gldim}R$ と置く。
$D$ の相対対角解消の長さは $2N$ であり、
$\operatorname{gldim}D\leq2N+d$ である。
有限生成 $C$-加群は有限生成 $D$-加群でもあるから、
有限生成二重次数射影 $D$-加群による有限解消を持つ。
この解消に二つの ambient saturation を適用する。
ambient 線束の cohomology は各因子で次数零と $N-1$ にのみあり、
各内部次数の群は有限射影 $R$-加群である。
従って線束から coherent 対象への Hom は有界で有限 $R$-cohomology を持つ。
$R$ の正則性から、これは $R$-perfect である。
次数別 $R$-flat な対象の解消では全 syzygy が次数別 $R$-flat なので、
この計算は係数変更を保つ。

一般の coherent な $M$ には有限線束和による上に有界な解消を使う。
各全次数に寄与する hyper-Ext spectral sequence の項は有限個であり、
直前の有限性から各 $\operatorname{Ext}^q(M,Q)$ は有限となる。
同じ計算は中心的な局所化と可換である。
有限複体を始域または終域に置いてもこの局所化の主張は変わらない。
\end{proof}

\subsubsection{正則局所底における compact 性の持上げ}

\begin{lemma}[通常の平坦対象の有限線束複体への retract]
\label{iff:lem:regular-base-compact-lift}
$(R,\mathfrak m,k)$ を上の族から得られる正則局所環とし、
$\mathcal C_R$ を一変数または二変数の射影圏とする。
coherent な $M\in\mathcal C_R$ が次数別 $R$-flat な有限生成 graded 加群で表示され、
閉ファイバー $M_k$ が $D(\mathcal C_k)$ で compact であるとする。
このとき $M$ は有限線束和からなる有界複体の直和因子であり、compact である。
\end{lemma}

\begin{proof}
$d=\dim R$ に関する帰納法を行う。
$d=0$ では仮定の compact 性から下と同じ一様 Ext 上界を得て、
証明末尾の線束解消の切断を行えばよい。
$d>0$ のとき正則パラメータ $t\in\mathfrak m$ を選ぶ。
$\bar R=R/(t)$ と $\bar M=M\otimes_R\bar R$ は同じ仮定を満たすので、
帰納法により $\bar M$ は compact である。
補題\ref{iff:lem:projective-product}の線束の一様な cohomological bound と、
compact 対象が有限個の線束の有限な cone と直和因子から得られることにより、
ある $b$ があって全ての coherent な $\bar Q$ について
\begin{equation}\label{iff:eq:uniform-ext-bound}
 \operatorname{Ext}^q_{\mathcal C_{\bar R}}(\bar M,\bar Q)=0
                \qquad(q>b)
\end{equation}
となる。$b$ は $\bar Q$ に依存しない。

まず $Q$ も次数別 $R$-flat な加群で表示されるとする。
中心的な係数商 $i$ について、graded 加群上の導来 tensor--restriction 随伴は
torsion を保つので射影圏へ降りる。
$M$ の flatness から
\[
 R\operatorname{Hom}_{\mathcal C_R}(M,i_*\bar Q)
   =R\operatorname{Hom}_{\mathcal C_{\bar R}}(\bar M,\bar Q)
\]
を得る。$0\to Q\xrightarrow{t}Q\to i_*\bar Q\to0$ と
式\eqref{iff:eq:uniform-ext-bound}により、$q>b$ のとき
$t:\operatorname{Ext}^q(M,Q)\to\operatorname{Ext}^q(M,Q)$ は全射である。
この Ext は有限 $R$-加群であり、$t\in\mathfrak m$ なので、
Nakayama により零である。

一般の coherent な $Q$ に対しては、その有限生成 graded 代表を選び、
有限 graded free 加群による解消の最初の $d$ 段を取る。
各 free 項は次数別 $R$-flat である。
$\operatorname{gldim}R=d$ と Tor の dimension shifting により、
第 $d$ syzygy も次数別 $R$-flat となる。
これらを tails へ移し、直前の消滅と長完全列を順に使うと、
全ての coherent な $Q$ に対して $\operatorname{Ext}^q(M,Q)=0$ を $q>b$ で得る。
特に、$d>1$ の場合に free 加群の最初の核を flat と仮定してはいない。

今度は $M$ 自身に有限線束和から順に全射を取り、$r>b$ に対して
\[
 0\longrightarrow K_r\longrightarrow P_{r-1}\longrightarrow\cdots
 \longrightarrow P_0\longrightarrow M\longrightarrow0
\]
を作る。noetherian 性から $K_r$ は coherent である。
この実際の解消が定める境界射 $M\to K_r[r]$ は
$\operatorname{Ext}^r(M,K_r)=0$ により零である。
従って $M$ は $[P_{r-1}\to\cdots\to P_0]$ の直和因子である。
各線束は compact なので主張を得る。
これは特殊ファイバー上の cone や idempotent の形式的持上げではなく、
$R$ 上にある解消の境界射の消滅である。完備化や代数化を使わない。
\end{proof}

\subsubsection{実対角の spread と五次元型の比較核}

\begin{theorem}[分離点を含む相対 smooth proper な開集合]
\label{iff:thm:smooth-neighborhood}
式\eqref{iff:eq:families}の各族に対して、$o$ を含む Zariski 開集合
$U\subset\operatorname{Spec}R_0$ が存在し、$\Sigma_U$ と $\Gamma_U$ は
相対 smooth proper である。特に
\[
 \Sigma_U\in\operatorname{Perf}
  (\Sigma_U^{\mathrm{op}}\otimes_U^{\mathbf L}\Sigma_U),\qquad
 \Gamma_U\in\operatorname{Perf}
  (\Gamma_U^{\mathrm{op}}\otimes_U^{\mathbf L}\Gamma_U)
\]
という実際の相対 perfect 対角を持つ。
各複素点 $s\in U$ では線束模型は全 coherent 圏に一致し、
残余模型は全 coherent Kuznetsov 成分を与える。

必要なら $U$ をさらに縮めることで、
$\rho=L_{\mathcal O}\circ(1)$ は残余模型の相対自己同値となる。
これは $o$ における通常の rotation を持ち上げ、
Artin 環を含む底変更と両立する。
\end{theorem}

\begin{proof}
$R=(R_0)_{\mathfrak m_o}$ と置く。
通常の二重次数対角を
\[
 (\Delta_A^+)_{i,j}=A_{i+j}\quad(i,j\geq0),\qquad
 \mathcal D_A=\pi\Delta_A^+
\]
とする。左右乗法が $C=A^{\mathrm{op}}\otimes_R A$ の作用を与える。
$A$ は次数一で生成されるので、$\Delta_A^+$ は $(0,0)$ の元 $1$ で生成される。
従って $\mathcal D_A$ は通常の coherent 対象であり、
その表示加群は各二重次数で有限自由である。

$E^\vee=\bigoplus_{i=0}^{N-1}\mathcal O_{A^{\mathrm{op}}}(-i)$、
$F=E^\vee\boxtimes_RE$ とする。
正則な $R$ 上では ambient は有限大域次元を持ち、
中心正則三次による商 $A$ は左右有限入射次元を持つ。
有限生成 torsion 加群の次数濾過の商は、有限 $R$-加群 $T_j$ を
内部次数 $j$ に置いた augmentation 加群である。
$R$ の正則性と augmentation 双対から、その導来双対も
有限個の内部次数にのみ存在する。
従ってこの Gorenstein 双対は torsion を torsion に移し、
線束 $\mathcal O_A(i)$ を $\mathcal O_{A^{\mathrm{op}}}(-i)$ に移すため、
線束の厚生成圏へ降りる。
有限線束模型の双対、外積と対角の evaluation は写像
\begin{equation}\label{iff:eq:actual-diagonal}
 \Sigma_R^{\mathrm{op}}\otimes_R^{\mathbf L}\Sigma_R
       \longrightarrow R\operatorname{End}(F),\qquad
 \Sigma_R\longrightarrow R\operatorname{Hom}(F,\mathcal D_A)
\end{equation}
を与える。後者は両側作用を持つ写像である。
これらは graded tensor--Hom evaluation を二つの導来 saturation に通して作る。
補題\ref{iff:lem:projective-product}により各辺は $R$-perfect で、
閉ファイバーへの底変更を保つ。
閉ファイバーでは、左右の noetherian 性、augmentation 双対による $\chi$、
有限 torsion dimension を満たすため、
\cite[Proposition~3.42, Lemma~4.9, Corollary~4.12,
Proposition~4.14, Theorem~4.15, Remark~4.17]{CMOBallardFarman}
の対角計算が適用され、両写像は同型となる。
従って cone に対する導来 Nakayama から、式\eqref{iff:eq:actual-diagonal}は
$R$ 上でも同型である。
ここで文献の体上の定理を、仮定を確認せず相対基底に適用してはいない。
全外積線束は $F$ によって厚生成されるので、$F$ は射影積の compact generator である。

$m=0$ の閉ファイバーの残余圏は補題\ref{ebk:lem:boundary}により
滑らかな可換三次四様体の Kuznetsov 成分に実際に同値である。
$m=1$ では、$cu^3$ の Jacobian 環が有限次元なので、
主稿の定理\paperref{main}{csc:thm:comparison}を適用できる。
どちらも残余圏は smooth proper である。
例外列との接着の Hom は体上 perfect であり、残余圏の smoothness により
接着双加群も perfect なので、有限例外列を加えた全線束圏も smooth proper となる。
従って閉ファイバーの式\eqref{iff:eq:actual-diagonal}の下で
$\mathcal D_{A_o}$ は compact である。
補題\ref{iff:lem:regular-base-compact-lift}を $\mathcal D_A$ に適用すると、
有限線束複体 $P$ と
\[
 \mathcal D_A\xrightarrow{s}P\xrightarrow{r}\mathcal D_A,
                  \qquad rs=1
\]
を得る。これを compact generator $F$ で移せば、
$\Sigma_R$ の相対 perfect 対角を得る。

この構成を原点を含む実開集合へ移す。
より具体的には、$\Delta_A^+$ の有限 graded free 項による解消を最初から $R_0$ 上で選ぶ。
局所環上の証明が与える Ext 上界より大きい $r$ で切れば、
その境界類は
\[
 e\in\operatorname{Ext}^r_{\operatorname{QGr}^{(2)}C_{R_0}}
              (\mathcal D_A,\mathcal K_r),\qquad e_{\mathfrak m_o}=0
\]
となる。補題\ref{iff:lem:projective-product}の有限性と局所化の互換性から、
ある $f\in R_0\setminus\mathfrak m_o$ がこの類を消す。
従って $R_0[f^{-1}]$ 上で同じ有界線束複体への retract が存在する。
特殊ファイバー上の複体の微分や idempotent を別途持ち上げる必要はない。
式\eqref{iff:eq:actual-diagonal}の同型も、その proper な cone を消す同じ形の開集合上に広がる。
従って $U=D(f)$ 上に実際の有限複体と retract が存在し、
$\Sigma_U$ は相対 smooth となる。
有限例外列の admissible な直交成分に smooth proper 性が移るため、
$\Gamma_U$ も相対 smooth proper である。

各複素ファイバーで全 coherent 圏との一致を確認する。
coherent な $M$ に対し、補題\ref{iff:lem:projective-product}の一変数版から
$R\operatorname{Hom}(E,M)$ は体上 perfect である。
smooth な $\Sigma_s$ 上では、体上 perfect な dg 加群は
perfect な対角との tensor により $\Sigma_s$-perfect となる。
$E$ の compact 生成性から $M$ 自身が compact であり、
逆向きは有限線束複体が有界 coherent cohomology を持つことによる。
有限 mutation で残余成分にも同じ結論を得る。

最後に $L_{\mathcal O}\circ(1)$ は、例外列の半直交性から残余線束圏を保ち、
有限線束模型上の実関手として全係数上に定義される。
$U$ 上では、その graph 双加群は生成対象で評価すると底上 perfect であり、
smooth proper な $\Gamma_U^{\mathrm{op}}\otimes_U^{\mathbf L}\Gamma_U$ 上でも perfect となる。
従って相対双対による左右随伴が存在し、底変更を保つ。
原点での関手は既知の rotation 自己同値なので、単位と余単位の cone は原点で零である。
有限生成対象上で評価した cone は底上 perfect だから、
原点を含むさらに小さい Zariski 開で両者が零となる。
これで相対自己同値を得る。
\end{proof}

\begin{corollary}[五次元 Clifford 混合族の可換実現を持つ開集合]
\label{iff:cor:cubic-fivefold-open}
$m=1$、$c\ne0$ とする。
原点を含む Zariski 開集合 $V\subset\mathbb A^9_\Lambda$ が存在し、
各複素点の全 coherent Kuznetsov 成分は滑らかな可換三次五様体の
Kuznetsov 成分に dg Morita 同値となる。
$V$ 上では \'etale 局所的に滑らかな可換三次族と実際の perfect 比較核が存在する。
原点における分離型比較核を固定すれば、その核を原点の \'etale 近傍へ延長できる。
\end{corollary}

\begin{proof}
定理\ref{iff:thm:smooth-neighborhood}は、主稿の
定理\paperref{main}{grl:thm:open-locus}に必要な実相対 smooth proper 模型を与える。
原点には定理\paperref{main}{csc:thm:comparison}の比較核がある。
$n=5\ne4$ なので開性定理を適用し、原点を含む可換実現 locus を $V$ とすればよい。
\end{proof}

\begin{remark}[結論の範囲]
この開集合の構成は対角の有限解消と分母払いによる存在証明であり、
補集合の方程式や全 smooth locus は計算していない。
従って任意の $\Lambda$、または任意の滑らかな $\Lambda$ に対する可換比較ではない。
$m=0$ に対する本小節の結論は相対 smooth proper な九変数近傍と rotation の持上げである。
四次元を除外する開性定理はここへ適用せず、四次元での比較は
系\ref{ifk:cor:clifford-fourfold-neighborhood}の rotation を用いる議論で扱う。
比較核について rotation、指定周期、中心作用を全て同時に保存する主張も含めない。
通常の非被約底上で全 coherent 圏を線束模型と同一視することは、
有限型の dg 底変更模型を持つという主張とは別である。
\end{remark}

% !TEX root = ../../mixed-sklyanin-cubics.tex
% Current result: 2026-11-22 / v62.
\subsection{Clifford 九変数族の幾何的接類と局所比較}
\label{ifk:subsec:clifford-application}

中心変数を加えない九変数族は三次四様体の次元に属する。
本節では実相対模型と原点の rotation 付き核を用い、
主稿の rotation を持つ変形の形式認識を適用する。
まず残余圏の九つの接方向に非幾何的成分がないことを確認し、
次に原点近傍の実際の可換比較へ進む。

定理\ref{ifc:thm:family}の $m=0$ の族を取り、
$0\in\operatorname{Spec}R$ を $\Lambda=0$ とする。
原点の商と残余圏を
\[
 A_0=(S_1\otimes S_2)/(a_1\Omega_{1,0}+a_2\Omega_{2,0}),
 \qquad \mathcal K_0=\Ku_{\mathrm{coh}}(A_0)
\]
と書く。
各因子の点曲線の非零方程式を $f_i$ とし、
$X_0=V(a_1f_1+a_2f_2)\subset\mathbb P^5$ と置く。
補題\ref{ebk:lem:boundary}は、この $\mathcal K_0$ と $\Ku(X_0)$ の間に
実際の dg Morita 同値 $E$ と
\begin{equation}\label{ifk:eq:marked-origin}
 E\rho_0\simeq\rho_{X_0}E
\end{equation}
を、指定三乗周期とともに与える。
同補題は二捩れ因子を許すので $S_2=S(1,1,\tau)$ に適用できる。
この入力は単なる非標識 Morita 同値より強い。

\begin{proposition}[九方向の非幾何的成分の消滅]
\label{ifk:prop:clifford-geometric-ks}
定理\ref{iff:thm:smooth-neighborhood}の実相対模型に対し、
原点における残余圏の Kodaira--Spencer 写像は
\begin{equation}\label{ifk:eq:clifford-ks}
 \operatorname{Mat}_3(\C)
   \xrightarrow{\ \operatorname{KS}_{\mathcal K}\ }
 \mathrm{HH}^2(\mathcal K_0)
   \xrightarrow{\ E_*\ }
 \mathrm{HH}^2(\Ku(X_0))
\end{equation}
の像を $\operatorname{im}\eta_{X_0}\simeq(J_{a_1f_1+a_2f_2})_3$ に持つ。
特に、二つの非恒等 rotation 固有空間への射影は
九方向全てについて零である。
\end{proposition}

\begin{proof}
定理\ref{iff:thm:smooth-neighborhood}により、
原点を含む開集合上に実相対 smooth proper 模型と
内部 shift から来る自己同値 $\rho$ がある。
その Artin 底変更により、任意の
$\Lambda_{ab}=\epsilon M_{ab}$ に対して
$\mathcal K_0$ の実際の Morita 変形と $\rho_0$ の持上げを得る。
任意の $M$ に補題\paperref{main}{ifk:lem:lifted-autoequivalence}を適用すると、
$\operatorname{KS}_{\mathcal K}(M)$ は $\rho_0$ に固定される。
式\eqref{ifk:eq:marked-origin}は、実際の可逆双加群による共役作用を
intertwine するので、その $E_*$ 像は $\rho_{X_0}$ に固定される。
命題\paperref{main}{ifk:prop:geometric-fixed-part}から結論を得る。
\end{proof}

\begin{corollary}[四次元九変数族の原点近傍における可換比較]
\label{ifk:cor:clifford-fourfold-neighborhood}
$m=0$ の九変数族について、原点を含む Zariski 開集合
$V\subset\mathbb A^9_\Lambda$ が存在し、
各複素点の全 coherent Kuznetsov 成分は、
滑らかな可換三次四様体の Kuznetsov 成分に dg Morita 同値である。

具体的には、式\eqref{ifk:eq:marked-origin}の実際の核 $E$ を固定すると、
原点の \'etale 近傍 $V'\to\mathbb A^9_\Lambda$ 上に、
$X_0$ を延長する滑らかな三次四様体族と、
$E$ を延長する perfect な実 Morita 比較核を構成できる。
この同値は任意の導来底変更に保たれる。
\end{corollary}

\begin{proof}
定理\ref{iff:thm:smooth-neighborhood}は、原点を含む実開集合上に
相対 smooth proper な残余圏族と、その実際の相対自己同値 $\rho$ を与える。
補題\ref{ebk:lem:boundary}の核 $E$ は
式\eqref{ifk:eq:marked-origin}を満たす。
従って主稿の定理\paperref{main}{ifk:thm:rotation-neighborhood}の
全ての仮定が満たされ、その結論を適用できる。
複素ファイバーで残余模型が全 coherent 残余圏に一致することも、
定理\ref{iff:thm:smooth-neighborhood}の結論である。
\end{proof}

\begin{remark}[接類と局所比較の範囲]
\label{ifk:rem:tangent-scope}
命題\ref{ifc:prop:tangent}は周囲代数の
$\mathrm{HH}^2(S_1\otimes S_2)_{(-2,2)}$ における九方向の独立性を証明する。
命題\ref{ifk:prop:clifford-geometric-ks}は残余圏への像の位置を制限するものであり、
写像\eqref{ifk:eq:clifford-ks}の単射性や rank 九を証明するものではない。
実相対残余模型は代数の次数別自由性だけから得たのではなく、
定理\ref{iff:thm:smooth-neighborhood}の相対 perfect 対角と
実自己同値の構成を用いている。

系\ref{ifk:cor:clifford-fourfold-neighborhood}では、一次成分の消滅だけを
全次数の比較へ置き換えたのではない。
主稿の定理\paperref{main}{ifk:thm:rotation-formal}による
有限位数作用の固定形式部分の認識と、
比較核スタックの有限表示を合わせて、実際の近傍の核を得ている。
全ての smooth 係数への延長、境界での比較核の保持、
および核による全 rotation・周期・中心作用の同時保存は、依然として別の課題である。
\end{remark}

